% ==============================================================================
% BANCO DE PREGUNTAS: ACADEMIA NACIONAL DE SEGURIDAD PÚBLICA (ANSP)
% PRUEBA CULTURAL - NIVEL BÁSICO (250 PREGUNTAS)
% ==============================================================================
\documentclass[10pt,a4paper]{article}
\usepackage[utf8]{inputenc}
\usepackage[T1]{fontenc}
\usepackage[spanish,es-nodecimaldot]{babel}
\usepackage{amsmath,amssymb,amsfonts}
\usepackage{geometry}
\usepackage{multicol}
\usepackage{enumitem}
\usepackage{tcolorbox}
\usepackage{fancyhdr}
\usepackage{hyperref}

\geometry{top=2cm, bottom=2cm, left=1.8cm, right=1.8cm}
\pagestyle{fancy}
\fancyhf{}
\rhead{\small\scshape ANSP - Prueba Cultural Nivel Básico}
\lhead{\small\bfseries Banco de 250 Preguntas}
\cfoot{\thepage}
\renewcommand{\headrulewidth}{0.4pt}

\begin{document}
\begin{titlepage}
\centering
{\scshape\LARGE ACADEMIA NACIONAL DE SEGURIDAD PÚBLICA \par}
\vspace{0.8cm}
{\Large\bfseries CONTENIDOS REVISADOS 2025 \par}
\vspace{1cm}
{\huge\bfseries CUESTIONARIO GENERAL PRUEBA CULTURAL \par}
{\Large\bfseries NIVEL BÁSICO (250 PREGUNTAS DE OPCIÓN MÚLTIPLE) \par}
\vspace{1.5cm}
\begin{tcolorbox}[colback=blue!5!white,colframe=blue!60!black,title=\textbf{TEMARIO OFICIAL}]
\begin{enumerate}[label=\textbf{\arabic*.}]
  \item \textbf{Historia de El Salvador} (50 Preguntas)
  \item \textbf{Lenguaje} (12 Preguntas)
  \item \textbf{Literatura} (38 Preguntas)
  \item \textbf{Geografía de El Salvador y Centroamérica} (50 Preguntas)
  \item \textbf{Realidad Económica y Social de El Salvador} (9 Preguntas)
  \item \textbf{Constitución de la República y Derechos Humanos} (27 Preguntas)
  \item \textbf{Legislación Salvadoreña} (19 Preguntas)
  \item \textbf{Matemáticas} (22 Preguntas)
  \item \textbf{Estadística y Probabilidad} (13 Preguntas)
  \item \textbf{Ofimática e Inteligencia Artificial} (10 Preguntas)
\end{enumerate}
\end{tcolorbox}
\vfill
{\large San Salvador, El Salvador \par}
{\large \today \par}
\end{titlepage}
\newpage

% ==============================================================================
% MATERIA 1: HISTORIA DE EL SALVADOR
% ==============================================================================
\section{1. Historia de El Salvador (50 Preguntas)}

\begin{enumerate}[label=\textbf{\arabic*.}]
\setcounter{enumi}{0}
\item ¿En qué siglo llegaron los conquistadores españoles al actual territorio de El Salvador?
\begin{enumerate}[label=\textbf{\Alph*)}]
\begin{multicols}{2}
  \item Siglo XIV
  \item Siglo XV
  \item Siglo XVI
  \item Siglo XVII
\end{multicols}
\end{enumerate}
\vspace{0.15cm}
\item ¿En qué año se firmó el Acta de Independencia de Centroamérica de la corona española?
\begin{enumerate}[label=\textbf{\Alph*)}]
\begin{multicols}{2}
  \item 1811
  \item 1821
  \item 1824
  \item 1841
\end{multicols}
\end{enumerate}
\vspace{0.15cm}
\item ¿Quién lideró la expedición militar española de conquista en El Salvador en 1524?
\begin{enumerate}[label=\textbf{\Alph*)}]
\begin{multicols}{2}
  \item Hernán Cortés
  \item Pedro de Alvarado
  \item Cristóbal Colón
  \item Francisco Pizarro
\end{multicols}
\end{enumerate}
\vspace{0.15cm}
\item ¿Qué prócer salvadoreño es reconocido históricamente como el "Padre de la Patria"?
\begin{enumerate}[label=\textbf{\Alph*)}]
\begin{multicols}{2}
  \item José Simeón Cañas
  \item Manuel José Arce
  \item José Matías Delgado
  \item Pedro Pablo Castillo
\end{multicols}
\end{enumerate}
\vspace{0.15cm}
\item ¿Qué acontecimiento histórico lideró Anastasio Aquino en El Salvador en 1833?
\begin{enumerate}[label=\textbf{\Alph*)}]
\begin{multicols}{2}
  \item La firma de la primera constitución
  \item Una rebelión indígena y campesina de los Nonualcos
  \item El derrocamiento del gobierno colonial
  \item La fundación del ejército salvadoreño
\end{multicols}
\end{enumerate}
\vspace{0.15cm}
\item ¿En qué departamento salvadoreño se desarrolló principalmente la Insurrección de los Nonualcos liderada por Anastasio Aquino?
\begin{enumerate}[label=\textbf{\Alph*)}]
\begin{multicols}{2}
  \item Santa Ana
  \item La Paz
  \item San Miguel
  \item Chalatenango
\end{multicols}
\end{enumerate}
\vspace{0.15cm}
\item ¿Quién fue el primer presidente de la República electo democráticamente tras la firma de los Acuerdos de Paz en 1994?
\begin{enumerate}[label=\textbf{\Alph*)}]
\begin{multicols}{2}
  \item Alfredo Cristiani
  \item Armando Calderón Sol
  \item Francisco Flores
  \item Mauricio Funes
\end{multicols}
\end{enumerate}
\vspace{0.15cm}
\item ¿Contra qué país vecino libró El Salvador el conflicto armado conocido como la "Guerra de las 100 Horas"?
\begin{enumerate}[label=\textbf{\Alph*)}]
\begin{multicols}{2}
  \item Guatemala
  \item Nicaragua
  \item Honduras
  \item Costa Rica
\end{multicols}
\end{enumerate}
\vspace{0.15cm}
\item ¿Qué factores principales desencadenaron la Guerra de las 100 Horas entre El Salvador y Honduras en 1969?
\begin{enumerate}[label=\textbf{\Alph*)}]
\begin{multicols}{2}
  \item Disputas por el canal interoceánico
  \item Tensiones migratorias y agrarias exacerbadas por partidos de fútbol eliminatorios
  \item Invasión de tropas estadounidenses
  \item Diferendos ideológicos de la Guerra Fría
\end{multicols}
\end{enumerate}
\vspace{0.15cm}
\item ¿Cuántos años duró formalmente la Guerra Civil salvadoreña (1980–1992)?
\begin{enumerate}[label=\textbf{\Alph*)}]
\begin{multicols}{2}
  \item 8 años
  \item 10 años
  \item 12 años
  \item 15 años
\end{multicols}
\end{enumerate}
\vspace{0.15cm}
\item ¿Qué trágico acontecimiento en 1980 marcó el inicio y la radicalización del Conflicto Armado salvadoreño?
\begin{enumerate}[label=\textbf{\Alph*)}]
\begin{multicols}{2}
  \item La caída de la Junta de Gobierno
  \item El asesinato de Monseñor Óscar Arnulfo Romero
  \item La ofensiva guerrillera 'Hasta el Tope'
  \item El golpe de Estado de 1979
\end{multicols}
\end{enumerate}
\vspace{0.15cm}
\item ¿Cuál fue la principal coalición insurgente que combatió a las fuerzas gubernamentales durante el conflicto armado salvadoreño?
\begin{enumerate}[label=\textbf{\Alph*)}]
\begin{multicols}{2}
  \item FSLN
  \item FMLN
  \item URNG
  \item Contras
\end{multicols}
\end{enumerate}
\vspace{0.15cm}
\item ¿Cuál de las siguientes organizaciones paramilitares gubernamentales NO formó parte del FMLN?
\begin{enumerate}[label=\textbf{\Alph*)}]
\begin{multicols}{2}
  \item FPL (Fuerzas Populares de Liberación)
  \item ERP (Ejército Revolucionario del Pueblo)
  \item ORDEN (Organización Democrática Nacionalista)
  \item RN (Resistencia Nacional)
\end{multicols}
\end{enumerate}
\vspace{0.15cm}
\item ¿En qué año se firmaron los históricos Acuerdos de Paz que pusieron fin a la guerra civil salvadoreña?
\begin{enumerate}[label=\textbf{\Alph*)}]
\begin{multicols}{2}
  \item 1989
  \item 1990
  \item 1992
  \item 1994
\end{multicols}
\end{enumerate}
\vspace{0.15cm}
\item ¿En qué ciudad se firmaron solemnemente los Acuerdos de Paz de El Salvador el 16 de enero de 1992?
\begin{enumerate}[label=\textbf{\Alph*)}]
\begin{multicols}{2}
  \item San José, Costa Rica
  \item Ginebra, Suiza
  \item Ciudad de México (Castillo de Chapultepec)
  \item Nueva York, EE. UU.
\end{multicols}
\end{enumerate}
\vspace{0.15cm}
\item ¿Quién era el presidente constitucional de El Salvador durante la firma de los Acuerdos de Paz en 1992?
\begin{enumerate}[label=\textbf{\Alph*)}]
\begin{multicols}{2}
  \item José Napoleón Duarte
  \item Alfredo Cristiani Burkard
  \item Armando Calderón Sol
  \item Álvaro Magaña
\end{multicols}
\end{enumerate}
\vspace{0.15cm}
\item En el marco de los Acuerdos de Paz de 1992, ¿qué organismo se creó para investigar los más graves crímenes y violaciones a derechos humanos?
\begin{enumerate}[label=\textbf{\Alph*)}]
\begin{multicols}{2}
  \item Corte Penal Internacional
  \item Comisión de la Verdad para El Salvador
  \item Tribunal Supremo de Justicia
  \item Ministerio de Seguridad Pública
\end{multicols}
\end{enumerate}
\vspace{0.15cm}
\item ¿Qué fenómeno delictivo organizado surgió e impactó fuertemente la seguridad salvadoreña en la época de postguerra?
\begin{enumerate}[label=\textbf{\Alph*)}]
\begin{multicols}{2}
  \item Bandolerismo rural decimonónico
  \item Pandillas o maras transnacionales
  \item Cárteles de piratería marítima
  \item Contrabando colonial
\end{multicols}
\end{enumerate}
\vspace{0.15cm}
\item ¿Qué instituto político alcanzó la presidencia de la República por primera vez en las elecciones de 2009?
\begin{enumerate}[label=\textbf{\Alph*)}]
\begin{multicols}{2}
  \item ARENA
  \item PDC
  \item FMLN
  \item GANA
\end{multicols}
\end{enumerate}
\vspace{0.15cm}
\item ¿Quién ejerce la presidencia de la República de El Salvador para el período 2024–2029?
\begin{enumerate}[label=\textbf{\Alph*)}]
\begin{multicols}{2}
  \item Salvador Sánchez Cerén
  \item Nayib Armando Bukele Ortez
  \item Félix Ulloa
  \item Ernesto Castro
\end{multicols}
\end{enumerate}
\vspace{0.15cm}
\item ¿Qué mandatario salvadoreño implementó la estrategia de seguridad denominada "Plan Control Territorial"?
\begin{enumerate}[label=\textbf{\Alph*)}]
\begin{multicols}{2}
  \item Antonio Saca
  \item Mauricio Funes
  \item Salvador Sánchez Cerén
  \item Nayib Bukele
\end{multicols}
\end{enumerate}
\vspace{0.15cm}
\item Tras la histórica reducción de los índices de homicidios, ¿qué área es señalada como la principal prioridad nacional para El Salvador?
\begin{enumerate}[label=\textbf{\Alph*)}]
\begin{multicols}{2}
  \item El problema económico (costo de vida y empleo)
  \item La demarcación fronteriza
  \item El reclutamiento militar obligatorio
  \item La exploración aeroespacial
\end{multicols}
\end{enumerate}
\vspace{0.15cm}
\item ¿Qué pueblo originario presentó la más férrea y prolongada resistencia contra la invasión de Pedro de Alvarado?
\begin{enumerate}[label=\textbf{\Alph*)}]
\begin{multicols}{2}
  \item Incas
  \item Pipiles (Cuzcatlecos)
  \item Toltecas de México
  \item Guaraníes
\end{multicols}
\end{enumerate}
\vspace{0.15cm}
\item ¿En qué célebre batalla de 1524 los guerreros pipiles hirieron gravemente en una pierna a Pedro de Alvarado?
\begin{enumerate}[label=\textbf{\Alph*)}]
\begin{multicols}{2}
  \item Batalla de Otumba
  \item Batalla de Acajutla
  \item Batalla de Grijalva
  \item Batalla del Golfo
\end{multicols}
\end{enumerate}
\vspace{0.15cm}
\item ¿Qué líder indígena fue capturado en San Vicente y ejecutado por el gobierno en 1833 tras encabezar la rebelión de los Nonualcos?
\begin{enumerate}[label=\textbf{\Alph*)}]
\begin{multicols}{2}
  \item Atlacatl
  \item Atonal
  \item Anastasio Aquino
  \item Feliciano Ama
\end{multicols}
\end{enumerate}
\vspace{0.15cm}
\item ¿En qué año se decretó la abolición total de la esclavitud en las Provincias Unidas de Centroamérica?
\begin{enumerate}[label=\textbf{\Alph*)}]
\begin{multicols}{2}
  \item 1811
  \item 1821
  \item 1824
  \item 1850
\end{multicols}
\end{enumerate}
\vspace{0.15cm}
\item ¿Quién fue investido como el primer obispo de San Salvador tras la independencia?
\begin{enumerate}[label=\textbf{\Alph*)}]
\begin{multicols}{2}
  \item Óscar Arnulfo Romero
  \item José Matías Delgado
  \item Luis Chávez y González
  \item Arturo Rivera Damas
\end{multicols}
\end{enumerate}
\vspace{0.15cm}
\item ¿Qué general y presidente provisional salvadoreño fue derrocado en octubre de 1944 tras haber asumido luego de la renuncia de Maximiliano Hernández Martínez?
\begin{enumerate}[label=\textbf{\Alph*)}]
\begin{multicols}{2}
  \item Andrés Ignacio Menéndez
  \item Osmín Aguirre y Salinas
  \item Salvador Castaneda Castro
  \item Julio Adalberto Rivera
\end{multicols}
\end{enumerate}
\vspace{0.15cm}
\item ¿En qué año ocurrió el trágico levantamiento campesino e indígena y su posterior masacre estatal conocida como "La Matanza"?
\begin{enumerate}[label=\textbf{\Alph*)}]
\begin{multicols}{2}
  \item 1922
  \item 1932
  \item 1944
  \item 1969
\end{multicols}
\end{enumerate}
\vspace{0.15cm}
\item ¿Qué mandatario salvadoreño inauguró tramos fundamentales de la Carretera Panamericana en el país durante su gobierno (1956–1960)?
\begin{enumerate}[label=\textbf{\Alph*)}]
\begin{multicols}{2}
  \item Fidel Sánchez Hernández
  \item José María Lemus
  \item Arturo Armando Molina
  \item Carlos Humberto Romero
\end{multicols}
\end{enumerate}
\vspace{0.15cm}
\item ¿Qué potencia extranjera proveyó mayor financiamiento, asesoría y armamento al gobierno salvadoreño durante el conflicto armado?
\begin{enumerate}[label=\textbf{\Alph*)}]
\begin{multicols}{2}
  \item Unión Soviética
  \item Estados Unidos
  \item Francia
  \item China
\end{multicols}
\end{enumerate}
\vspace{0.15cm}
\item ¿Qué país caribeño brindó apoyo político, logístico y entrenamiento doctrinario a las fuerzas guerrilleras del FMLN?
\begin{enumerate}[label=\textbf{\Alph*)}]
\begin{multicols}{2}
  \item Jamaica
  \item República Dominicana
  \item Cuba
  \item Puerto Rico
\end{multicols}
\end{enumerate}
\vspace{0.15cm}
\item ¿Qué porcentaje aproximado de la población total salvadoreña de la época se estima que pereció a causa de la Guerra Civil?
\begin{enumerate}[label=\textbf{\Alph*)}]
\begin{multicols}{2}
  \item Alrededor del 2\% (aprox. 75,000 vidas)
  \item Cerca del 10\%
  \item Menos del 0.1\%
  \item Aproximadamente el 15\%
\end{multicols}
\end{enumerate}
\vspace{0.15cm}
\item ¿Qué hecho y lugar solemne simbolizó la culminación de las negociaciones de paz en 1992?
\begin{enumerate}[label=\textbf{\Alph*)}]
\begin{multicols}{2}
  \item La firma en el Castillo de Chapultepec
  \item La reunión en Ginebra
  \item El pacto de La Palma
  \item La cumbre de Esquipulas
\end{multicols}
\end{enumerate}
\vspace{0.15cm}
\item ¿Qué sector económico creció exponencialmente en El Salvador tras la firma de la paz, transformando el modelo agrario tradicional?
\begin{enumerate}[label=\textbf{\Alph*)}]
\begin{multicols}{2}
  \item Sector minero pesado
  \item Sector servicios, comercio y flujo de remesas familiares
  \item Exportación algodonera
  \item Industria aeroespacial
\end{multicols}
\end{enumerate}
\vspace{0.15cm}
\item ¿En qué año entró en vigencia la Ley de Integración Monetaria que adoptó el dólar estadounidense como moneda de curso legal?
\begin{enumerate}[label=\textbf{\Alph*)}]
\begin{multicols}{2}
  \item 1992
  \item 1998
  \item 2001
  \item 2009
\end{multicols}
\end{enumerate}
\vspace{0.15cm}
\item ¿Qué devastador huracán golpeó a El Salvador y Centroamérica a finales de 1998 provocando graves inundaciones y deslaves?
\begin{enumerate}[label=\textbf{\Alph*)}]
\begin{multicols}{2}
  \item Huracán Katrina
  \item Huracán Mitch
  \item Huracán Fifi
  \item Tormenta Stan
\end{multicols}
\end{enumerate}
\vspace{0.15cm}
\item ¿En qué año ocurrieron dos devastadores terremotos de gran magnitud (enero y febrero) que destruyeron miles de viviendas salvadoreñas?
\begin{enumerate}[label=\textbf{\Alph*)}]
\begin{multicols}{2}
  \item 1986
  \item 1998
  \item 2001
  \item 2011
\end{multicols}
\end{enumerate}
\vspace{0.15cm}
\item ¿En qué año aprobó la Asamblea Legislativa salvadoreña el Régimen de Excepción para combatir a las estructuras terroristas y pandillas?
\begin{enumerate}[label=\textbf{\Alph*)}]
\begin{multicols}{2}
  \item 2019
  \item 2020
  \item 2021
  \item 2022
\end{multicols}
\end{enumerate}
\vspace{0.15cm}
\item ¿Qué presidente centroamericano recibió el Premio Nobel de la Paz en 1987 por impulsar los Acuerdos de Esquipulas para pacificar El Salvador y la región?
\begin{enumerate}[label=\textbf{\Alph*)}]
\begin{multicols}{2}
  \item José Napoleón Duarte
  \item Óscar Arias Sánchez (Costa Rica)
  \item Vinicio Cerezo
  \item Daniel Ortega
\end{multicols}
\end{enumerate}
\vspace{0.15cm}
\item ¿Qué gran arteria vial de conectividad interdepartamental fue inaugurada en 2012 con fondos de la Corporación del Reto del Milenio?
\begin{enumerate}[label=\textbf{\Alph*)}]
\begin{multicols}{2}
  \item Carretera del Litoral
  \item Carretera Longitudinal del Norte
  \item Bulevar Monseñor Romero
  \item Autopista a Comalapa
\end{multicols}
\end{enumerate}
\vspace{0.15cm}
\item ¿En qué año proclamó El Salvador su separación definitiva y soberana de la República Federal de Centroamérica?
\begin{enumerate}[label=\textbf{\Alph*)}]
\begin{multicols}{2}
  \item 1821
  \item 1838
  \item 1841
  \item 1856
\end{multicols}
\end{enumerate}
\vspace{0.15cm}
\item ¿Qué instrumento jurídico internacional definió los límites territoriales y áreas de controversia ("bolsones") entre El Salvador y Honduras en 1992?
\begin{enumerate}[label=\textbf{\Alph*)}]
\begin{multicols}{2}
  \item Tratado de Paz de Lima
  \item Fallo de la Corte Internacional de Justicia de La Haya (Tratado de Límites)
  \item Pacto de San Salvador
  \item Convenio de Tegucigalpa
\end{multicols}
\end{enumerate}
\vspace{0.15cm}
\item ¿Qué efervescencia sociopolítica caracterizó el panorama salvadoreño durante la década de 1970 previo a la guerra civil?
\begin{enumerate}[label=\textbf{\Alph*)}]
\begin{multicols}{2}
  \item Surgimiento masivo de organizaciones populares, gremiales y células político-militares
  \item Periodo de completa estabilidad y desmilitarización
  \item Abolición del ejército
  \item Guerra con Guatemala
\end{multicols}
\end{enumerate}
\vspace{0.15cm}
\item ¿Qué emblemática central hidroeléctrica fue inaugurada en 1954 sobre el río Lempa bajo la presidencia del Tte. Cnel. Óscar Osorio?
\begin{enumerate}[label=\textbf{\Alph*)}]
\begin{multicols}{2}
  \item Presa Cerrón Grande
  \item Central Hidroeléctrica 5 de Noviembre (Chorrera del Guayabo)
  \item Presa 15 de Septiembre
  \item Central 3 de Febrero (El Chaparral)
\end{multicols}
\end{enumerate}
\vspace{0.15cm}
\item ¿En qué año se fundó la Universidad de El Salvador (UES), el centro de educación superior más antiguo del país?
\begin{enumerate}[label=\textbf{\Alph*)}]
\begin{multicols}{2}
  \item 1821
  \item 1841
  \item 1886
  \item 1905
\end{multicols}
\end{enumerate}
\vspace{0.15cm}
\item ¿Qué trascendental institución de seguridad social y salud pública fue creada durante la gestión del presidente Óscar Osorio (1950–1956)?
\begin{enumerate}[label=\textbf{\Alph*)}]
\begin{multicols}{2}
  \item Ministerio de Hacienda
  \item Instituto Salvadoreño del Seguro Social (ISSS)
  \item Banco Central de Reserva
  \item Fondo Social para la Vivienda
\end{multicols}
\end{enumerate}
\vspace{0.15cm}
\item ¿En qué siglo se consolidó el cultivo del café como el principal producto agroexportador de la economía salvadoreña?
\begin{enumerate}[label=\textbf{\Alph*)}]
\begin{multicols}{2}
  \item Siglo XVII
  \item Siglo XVIII
  \item Siglo XIX
  \item Siglo XXI
\end{multicols}
\end{enumerate}
\vspace{0.15cm}
\item ¿Qué gobierno decretó la nacionalización de la banca y la reforma agraria en marzo de 1980?
\begin{enumerate}[label=\textbf{\Alph*)}]
\begin{multicols}{2}
  \item Maximiliano Hernández Martínez
  \item Junta Revolucionaria de Gobierno
  \item Alfredo Cristiani
  \item Armando Calderón Sol
\end{multicols}
\end{enumerate}
\vspace{0.15cm}
\item ¿Qué gesta popular y patriótica es reconocida como el "Primer Grito de Independencia" de Centroamérica?
\begin{enumerate}[label=\textbf{\Alph*)}]
\begin{multicols}{2}
  \item El alzamiento del 5 de noviembre de 1811 en San Salvador
  \item La batalla de San Jacinto en 1856
  \item La conspiración de Belén en 1813
  \item La toma de Granada en 1823
\end{multicols}
\end{enumerate}
\vspace{0.15cm}
\end{enumerate}\n
\vspace{0.4cm}
\begin{tcolorbox}[colback=gray!5!white,colframe=gray!60!black,title=\textbf{CLAVE DE RESPUESTAS -- HISTORIA DE EL SALVADOR}]
\begin{multicols}{5}
\begin{enumerate}[label=\textbf{\arabic*.}]
\setcounter{enumi}{0}
  \item \textbf{C}
  \item \textbf{B}
  \item \textbf{B}
  \item \textbf{C}
  \item \textbf{B}
  \item \textbf{B}
  \item \textbf{B}
  \item \textbf{C}
  \item \textbf{B}
  \item \textbf{C}
  \item \textbf{B}
  \item \textbf{B}
  \item \textbf{C}
  \item \textbf{C}
  \item \textbf{C}
  \item \textbf{B}
  \item \textbf{B}
  \item \textbf{B}
  \item \textbf{C}
  \item \textbf{B}
  \item \textbf{D}
  \item \textbf{A}
  \item \textbf{B}
  \item \textbf{B}
  \item \textbf{C}
  \item \textbf{C}
  \item \textbf{B}
  \item \textbf{A}
  \item \textbf{B}
  \item \textbf{B}
  \item \textbf{B}
  \item \textbf{C}
  \item \textbf{A}
  \item \textbf{A}
  \item \textbf{B}
  \item \textbf{C}
  \item \textbf{B}
  \item \textbf{C}
  \item \textbf{D}
  \item \textbf{B}
  \item \textbf{B}
  \item \textbf{C}
  \item \textbf{B}
  \item \textbf{A}
  \item \textbf{B}
  \item \textbf{B}
  \item \textbf{B}
  \item \textbf{C}
  \item \textbf{B}
  \item \textbf{A}
\end{enumerate}
\end{multicols}
\end{tcolorbox}
\newpage\n
% ==============================================================================
% MATERIA 2: LENGUAJE
% ==============================================================================
\section{2. Lenguaje (11 Preguntas)}

\begin{enumerate}[label=\textbf{\arabic*.}]
\setcounter{enumi}{50}
\item Completa correctamente la oración: "El niño canta una __"
\begin{enumerate}[label=\textbf{\Alph*)}]
\begin{multicols}{2}
  \item canción
  \item cancion
  \item cansion
  \item canzión
\end{multicols}
\end{enumerate}
\vspace{0.15cm}
\item ¿Cuál de los siguientes enunciados presenta una redacción y concordancia gramatical correctas?
\begin{enumerate}[label=\textbf{\Alph*)}]
\begin{multicols}{2}
  \item Ayer fuistes al cine con tus amigos.
  \item Ayer fui al cine con mis amigos.
  \item Ayer iba al cine con los amigos míos ayer.
  \item Nosotros fui al cine ayer.
\end{multicols}
\end{enumerate}
\vspace{0.15cm}
\item ¿Cuál es la grafía correcta de la locución que expresa expectativa o comprobación?
\begin{enumerate}[label=\textbf{\Alph*)}]
\begin{multicols}{2}
  \item Haber si vamos
  \item A ver si vamos
  \item Ha ver si vamos
  \item Aver si vamos
\end{multicols}
\end{enumerate}
\vspace{0.15cm}
\item ¿Cuál de las siguientes formas verbales lleva tilde ortográfica correctamente?
\begin{enumerate}[label=\textbf{\Alph*)}]
\begin{multicols}{2}
  \item Esta (verbo)
  \item Está (del verbo estar)
  \item Ésta (verbo)
  \item Éstar
\end{multicols}
\end{enumerate}
\vspace{0.15cm}
\item En expresión escrita y estilística, ¿en qué consiste el error sintáctico conocido como "anacoluto"?
\begin{enumerate}[label=\textbf{\Alph*)}]
\begin{multicols}{2}
  \item Repetir una misma palabra en el mismo párrafo
  \item Falta de coherencia sintáctica por ruptura de la estructura oracional iniciada
  \item Confundir la 'b' con la 'v'
  \item Omitir signos de interrogación
\end{multicols}
\end{enumerate}
\vspace{0.15cm}
\item ¿Qué elemento estructural y lógico es indispensable y NO puede faltar en un ensayo argumentativo?
\begin{enumerate}[label=\textbf{\Alph*)}]
\begin{multicols}{2}
  \item Rima asonante
  \item Tesis o postura central que se defiende
  \item Diálogo dramatizado
  \item Uso obligatorio de metáforas
\end{multicols}
\end{enumerate}
\vspace{0.15cm}
\item ¿Qué error metodológico debe evitarse rigurosamente en la redacción de una conclusión?
\begin{enumerate}[label=\textbf{\Alph*)}]
\begin{multicols}{2}
  \item Sintetizar los puntos clave
  \item Introducir nuevos argumentos no expuestos previamente en el desarrollo
  \item Reiterar la postura del autor
  \item Plantear reflexiones finales
\end{multicols}
\end{enumerate}
\vspace{0.15cm}
\item ¿Cuál es un precepto elemental e inexcusable en el ejercicio de la oratoria efectiva?
\begin{enumerate}[label=\textbf{\Alph*)}]
\begin{multicols}{2}
  \item Hablar a gran velocidad sin pausas
  \item Hablar con claridad, articulación y dicción nítida
  \item Mirar fijamente hacia el suelo
  \item Usar solo tecnicismos incomprensibles
\end{multicols}
\end{enumerate}
\vspace{0.15cm}
\item ¿Qué actitud no verbal es esencial para generar confianza y empatía con el auditorio?
\begin{enumerate}[label=\textbf{\Alph*)}]
\begin{multicols}{2}
  \item Cruzar los brazos y mostrar enfado
  \item Lenguaje corporal abierto, contacto visual y postura erguida
  \item Permanecer de espaldas al público
  \item Leer continuamente notas sin alzar la mirada
\end{multicols}
\end{enumerate}
\vspace{0.15cm}
\item ¿Qué técnica psicofisiológica es altamente recomendada para regular la ansiedad y el miedo escénico al hablar en público?
\begin{enumerate}[label=\textbf{\Alph*)}]
\begin{multicols}{2}
  \item Hiperventilación acelerada
  \item Respiración diafragmática profunda y controlada
  \item Ingerir bebidas energéticas en exceso
  \item Hablar sin respirar
\end{multicols}
\end{enumerate}
\vspace{0.15cm}
\item ¿Cuál de las siguientes conductas NO constituye una cualidad de un orador eficaz y profesional?
\begin{enumerate}[label=\textbf{\Alph*)}]
\begin{multicols}{2}
  \item Adaptar el mensaje al perfil del público
  \item Uso excesivo y pedante de jerga técnica innecesaria
  \item Manejo equilibrado de tiempos y silencios
  \item Capacidad de escucha y modulación
\end{multicols}
\end{enumerate}
\vspace{0.15cm}
\end{enumerate}\n
\vspace{0.4cm}
\begin{tcolorbox}[colback=gray!5!white,colframe=gray!60!black,title=\textbf{CLAVE DE RESPUESTAS -- LENGUAJE}]
\begin{multicols}{5}
\begin{enumerate}[label=\textbf{\arabic*.}]
\setcounter{enumi}{50}
  \item \textbf{A}
  \item \textbf{B}
  \item \textbf{B}
  \item \textbf{B}
  \item \textbf{B}
  \item \textbf{B}
  \item \textbf{B}
  \item \textbf{B}
  \item \textbf{B}
  \item \textbf{B}
  \item \textbf{B}
\end{enumerate}
\end{multicols}
\end{tcolorbox}
\newpage\n
% ==============================================================================
% MATERIA 3: LITERATURA
% ==============================================================================
\section{3. Literatura (39 Preguntas)}

\begin{enumerate}[label=\textbf{\arabic*.}]
\setcounter{enumi}{57}
\item El célebre escritor salvadoreño Alfredo Espino es universalmente recordado por cultivar:
\begin{enumerate}[label=\textbf{\Alph*)}]
\begin{multicols}{2}
  \item Tratados de sociología
  \item Poemas líricos y bucólicos
  \item Novelas policiacas
  \item Ensayos de economía
\end{multicols}
\end{enumerate}
\vspace{0.15cm}
\item ¿Qué tipo de obra literaria es la insigne publicación "Jícaras Tristes" de Alfredo Espino?
\begin{enumerate}[label=\textbf{\Alph*)}]
\begin{multicols}{2}
  \item Pieza de teatro
  \item Poemario (libro de poesía lírica)
  \item Crónica periodística
  \item Novela histórica
\end{multicols}
\end{enumerate}
\vspace{0.15cm}
\item ¿Qué temática principal predomina en los versos del poemario "Jícaras Tristes"?
\begin{enumerate}[label=\textbf{\Alph*)}]
\begin{multicols}{2}
  \item La guerra nuclear y el cosmos
  \item La naturaleza, el paisaje rural y la identidad del campo salvadoreño
  \item La mitología griega clásica
  \item La industrialización de las urbes
\end{multicols}
\end{enumerate}
\vspace{0.15cm}
\item ¿Qué recurso retórico figura predominantemente en la imaginería poética de "Jícaras Tristes"?
\begin{enumerate}[label=\textbf{\Alph*)}]
\begin{multicols}{2}
  \item Metáfora y personificación de la naturaleza
  \item Jitanjáfora
  \item Oxímoron barroco
  \item Ironía sarcástica
\end{multicols}
\end{enumerate}
\vspace{0.15cm}
\item ¿Quién es el inmortal dramaturgo inglés autor de la tragedia "Romeo y Julieta"?
\begin{enumerate}[label=\textbf{\Alph*)}]
\begin{multicols}{2}
  \item Christopher Marlowe
  \item William Shakespeare
  \item John Milton
  \item Lord Byron
\end{multicols}
\end{enumerate}
\vspace{0.15cm}
\item ¿En qué ciudad italiana se desarrollan los trágicos sucesos de la obra "Romeo y Julieta"?
\begin{enumerate}[label=\textbf{\Alph*)}]
\begin{multicols}{2}
  \item Roma
  \item Florencia
  \item Verona
  \item Venecia
\end{multicols}
\end{enumerate}
\vspace{0.15cm}
\item ¿Cuál es la fuerza o causa antagonista primordial que condena al fracaso el amor de Romeo y Julieta?
\begin{enumerate}[label=\textbf{\Alph*)}]
\begin{multicols}{2}
  \item Una epidemia de peste
  \item El odio y la enemistad ancestral entre las familias Montesco y Capuleto
  \item La intervención militar de Mantua
  \item La traición de Fray Lorenzo
\end{multicols}
\end{enumerate}
\vspace{0.15cm}
\item ¿Qué laureado escritor colombiano y exponente del Realismo Mágico escribió "Cien años de soledad"?
\begin{enumerate}[label=\textbf{\Alph*)}]
\begin{multicols}{2}
  \item Mario Vargas Llosa
  \item Gabriel García Márquez
  \item Julio Cortázar
  \item Carlos Fuentes
\end{multicols}
\end{enumerate}
\vspace{0.15cm}
\item ¿Cómo se llama el mítico pueblo ficticio fundado por la estirpe de los Buendía en "Cien años de soledad"?
\begin{enumerate}[label=\textbf{\Alph*)}]
\begin{multicols}{2}
  \item Comala
  \item Macondo
  \item Santa María
  \item Yoknapatawpha
\end{multicols}
\end{enumerate}
\vspace{0.15cm}
\item En "Cien años de soledad", ¿qué significado simbólico entraña el episodio de la "peste del insomnio y del olvido"?
\begin{enumerate}[label=\textbf{\Alph*)}]
\begin{multicols}{2}
  \item La necesidad de medicinas modernas
  \item La pérdida de la memoria colectiva y la identidad histórica
  \item El progreso industrial de la comarca
  \item El castigo divino por la guerra
\end{multicols}
\end{enumerate}
\vspace{0.15cm}
\item ¿Qué prolífico artista y escritor salvadoreño es el autor del clásico libro de narrativa costumbrista "Cuentos de Barro"?
\begin{enumerate}[label=\textbf{\Alph*)}]
\begin{multicols}{2}
  \item Roque Dalton
  \item Salarrué (Salvador Salazar Arrué)
  \item Francisco Gavidia
  \item Arturo Ambrogi
\end{multicols}
\end{enumerate}
\vspace{0.15cm}
\item ¿Qué corriente de renovación estética universal influyó profundamente en los escritores y poetas salvadoreños de mediados del siglo XX?
\begin{enumerate}[label=\textbf{\Alph*)}]
\begin{multicols}{2}
  \item Neoclasicismo académico
  \item Vanguardismo literario y poesía comprometida
  \item Mester de Clerecía
  \item Romanticismo tardío decimonónico
\end{multicols}
\end{enumerate}
\vspace{0.15cm}
\item ¿Qué erudito y polígrafo es unánimemente considerado el padre fundador y pionero de la literatura salvadoreña moderna?
\begin{enumerate}[label=\textbf{\Alph*)}]
\begin{multicols}{2}
  \item Alberto Masferrer
  \item Francisco Gavidia
  \item David J. Guzmán
  \item Vicente Acosta
\end{multicols}
\end{enumerate}
\vspace{0.15cm}
\item ¿Quién escribió el clásico texto poético costumbrista sobre la naturaleza cuscatleca titulado "El libro del Trópico"?
\begin{enumerate}[label=\textbf{\Alph*)}]
\begin{multicols}{2}
  \item Claudia Lars
  \item Arturo Ambrogi
  \item Francisco Gavidia
  \item Alfredo Espino
\end{multicols}
\end{enumerate}
\vspace{0.15cm}
\item ¿Qué libro de Roque Dalton conjuga el humor, la ironía militante y la crítica política bajo seudónimos poéticos?
\begin{enumerate}[label=\textbf{\Alph*)}]
\begin{multicols}{2}
  \item Las historias prohibidas del Pulgarcito
  \item Poemas Clandestinos
  \item Taberna y otros lugares
  \item Pobrecito poeta que era yo
\end{multicols}
\end{enumerate}
\vspace{0.15cm}
\item ¿Qué novelista salvadoreño contemporáneo, autor de "Un día en la vida", fue galardonado con el Premio Nacional de Cultura en 2009?
\begin{enumerate}[label=\textbf{\Alph*)}]
\begin{multicols}{2}
  \item Horacio Castellanos Moya
  \item Manlio Argueta
  \item Rafael Menjívar Ochoa
  \item David Escobar Galindo
\end{multicols}
\end{enumerate}
\vspace{0.15cm}
\item ¿Qué aclamada novela de Manlio Argueta retrata el sufrimiento campesino y la represión militar previa y durante la guerra civil salvadoreña?
\begin{enumerate}[label=\textbf{\Alph*)}]
\begin{multicols}{2}
  \item Caperucita en la zona roja
  \item Un día en la vida
  \item Cuzcatlán donde bate la mar del sur
  \item El valle de las hamacas
\end{multicols}
\end{enumerate}
\vspace{0.15cm}
\item ¿Qué insigne poetisa salvadoreña es la autora de la colección lírica "Romances de Norte y Sur"?
\begin{enumerate}[label=\textbf{\Alph*)}]
\begin{multicols}{2}
  \item Prudencia Ayala
  \item Claudia Lars (Margarita del Carmen Brannon)
  \item Liliam Jiménez
  \item Claribel Alegría
\end{multicols}
\end{enumerate}
\vspace{0.15cm}
\item ¿Qué influyente publicación periódica literaria fundó Francisco Gavidia a finales del siglo XIX para difundir las ciencias y letras?
\begin{enumerate}[label=\textbf{\Alph*)}]
\begin{multicols}{2}
  \item La Hoja
  \item El Repertorio Salvadoreño
  \item Cultura
  \item El Eco de Occidente
\end{multicols}
\end{enumerate}
\vspace{0.15cm}
\item ¿Qué drama teatral escrito en verso por Francisco Gavidia aborda la temática indígena y el poder colonial?
\begin{enumerate}[label=\textbf{\Alph*)}]
\begin{multicols}{2}
  \item La princesa Citalá
  \item Ursino
  \item Lucía Jerez
  \item Ollantay
\end{multicols}
\end{enumerate}
\vspace{0.15cm}
\item ¿Qué distinguido escritor y abogado salvadoreño, hermano de Alfredo Espino, escribió "Mitología de Cuscatlán" y "Trenes"?
\begin{enumerate}[label=\textbf{\Alph*)}]
\begin{multicols}{2}
  \item Miguel Ángel Espino
  \item Salarrué
  \item Alberto Rivas Bonilla
  \item Trigueros de León
\end{multicols}
\end{enumerate}
\vspace{0.15cm}
\item ¿Qué prominente poeta salvadoreño integró en los años 50 la combativa "Generación Comprometida"?
\begin{enumerate}[label=\textbf{\Alph*)}]
\begin{multicols}{2}
  \item José Roberto Cea
  \item Roque Dalton García
  \item Ítalo López Vallecillos
  \item Todos los anteriores
\end{multicols}
\end{enumerate}
\vspace{0.15cm}
\item ¿Qué valiente escritora y pionera sufragista salvadoreña intentó postularse a la presidencia en 1930 desafiando las leyes de la época?
\begin{enumerate}[label=\textbf{\Alph*)}]
\begin{multicols}{2}
  \item Claudia Lars
  \item Prudencia Ayala
  \item Antonia Navarro
  \item Ana Dolores Arias
\end{multicols}
\end{enumerate}
\vspace{0.15cm}
\item ¿Qué figura retórica establece una comparación explícita entre dos términos mediante nexos como "como", "cual" o "semejante a"?
\begin{enumerate}[label=\textbf{\Alph*)}]
\begin{multicols}{2}
  \item Metáfora
  \item Símil o comparación
  \item Hipérbole
  \item Antítesis
\end{multicols}
\end{enumerate}
\vspace{0.15cm}
\item ¿Qué corriente o subgénero novelístico define la cosmovisión y técnica de "Cien años de soledad"?
\begin{enumerate}[label=\textbf{\Alph*)}]
\begin{multicols}{2}
  \item Novela gótica
  \item Novela del Realismo Mágico
  \item Ciencia ficción distópica
  \item Naturalismo experimental
\end{multicols}
\end{enumerate}
\vspace{0.15cm}
\item ¿Qué otra famosa tragedia dramática de William Shakespeare está ambientada en la península italiana?
\begin{enumerate}[label=\textbf{\Alph*)}]
\begin{multicols}{2}
  \item Hamlet
  \item Otelo, el moro de Venecia
  \item Macbeth
  \item El rey Lear
\end{multicols}
\end{enumerate}
\vspace{0.15cm}
\item ¿Qué excelso poeta nicaragüense y líder del Modernismo hispanoamericano mantuvo estrecha amistad e intercambio con Francisco Gavidia?
\begin{enumerate}[label=\textbf{\Alph*)}]
\begin{multicols}{2}
  \item Ernesto Cardenal
  \item Rubén Darío
  \item Alfonso Cortés
  \item Salomón de la Selva
\end{multicols}
\end{enumerate}
\vspace{0.15cm}
\item ¿Qué trascendental distinción internacional otorgó la Academia Sueca a Gabriel García Márquez en 1982?
\begin{enumerate}[label=\textbf{\Alph*)}]
\begin{multicols}{2}
  \item Premio Cervantes
  \item Premio Nobel de Literatura
  \item Premio Pulitzer
  \item Premio Rómulo Gallegos
\end{multicols}
\end{enumerate}
\vspace{0.15cm}
\item ¿Cuál es la característica medular del Realismo Mágico en la narrativa latinoamericana?
\begin{enumerate}[label=\textbf{\Alph*)}]
\begin{multicols}{2}
  \item El rigor científico y matemático
  \item La inclusión e integración de sucesos maravillosos o fantásticos en la vida cotidiana sin asombro de los personajes
  \item La denuncia policial sin ficción
  \item El rechazo a la mitología popular
\end{multicols}
\end{enumerate}
\vspace{0.15cm}
\item ¿Qué conocida obra teatral para niños y adultos escribió la poetisa salvadoreña Claudia Lars?
\begin{enumerate}[label=\textbf{\Alph*)}]
\begin{multicols}{2}
  \item La olla de barro
  \item El paraíso de los imprudentes (o Fábulas del patio)
  \item La casa de Bernarda Alba
  \item La boda del coyote
\end{multicols}
\end{enumerate}
\vspace{0.15cm}
\item ¿Qué conmovedor libro autobiográfico en prosa poética escribió Claudia Lars rememorando sus años en Armenia, Sonsonate?
\begin{enumerate}[label=\textbf{\Alph*)}]
\begin{multicols}{2}
  \item Estrellas en el pozo
  \item Tierra de infancia
  \item Donde llegan los pasos
  \item Canción redonda
\end{multicols}
\end{enumerate}
\vspace{0.15cm}
\item ¿En qué año fue trágicamente asesinado el revolucionario y poeta Roque Dalton?
\begin{enumerate}[label=\textbf{\Alph*)}]
\begin{multicols}{2}
  \item 1969
  \item 1975
  \item 1980
  \item 1992
\end{multicols}
\end{enumerate}
\vspace{0.15cm}
\item ¿Qué conocida antología poética de clásicos y autores centroamericanos fue editada y compilada bajo el título evocador de Frazer?
\begin{enumerate}[label=\textbf{\Alph*)}]
\begin{multicols}{2}
  \item Las mil y una noches
  \item La rama de oro
  \item El cancionero de Amberes
  \item El libro de las tierras vírgenes
\end{multicols}
\end{enumerate}
\vspace{0.15cm}
\item ¿Cuál es el medio y vehículo de comunicación literaria más antiguo de la humanidad?
\begin{enumerate}[label=\textbf{\Alph*)}]
\begin{multicols}{2}
  \item La imprenta de tipos móviles
  \item La oralidad (palabra hablada y cantada)
  \item El pergamino manuscrito
  \item Las tablillas cuneiformes
\end{multicols}
\end{enumerate}
\vspace{0.15cm}
\item ¿Qué rasgo distingue esencialmente a la literatura de tradición oral en su momento de difusión?
\begin{enumerate}[label=\textbf{\Alph*)}]
\begin{multicols}{2}
  \item Ser inmutable y fija en papel
  \item Su transmisión mediante la voz, la entonación y la puesta en escena (performance)
  \item La firma de derechos de autor registrados
  \item El uso exclusivo de prosa científica
\end{multicols}
\end{enumerate}
\vspace{0.15cm}
\item ¿En qué consiste el sistema de "escritura ideográfica" desarrollado por civilizaciones antiguas?
\begin{enumerate}[label=\textbf{\Alph*)}]
\begin{multicols}{2}
  \item En representar sonidos alfabéticos individuales
  \item En representar conceptos e ideas complejas a través de símbolos o figuras gráficas
  \item En grabar solo números
  \item En imitar el código Morse
\end{multicols}
\end{enumerate}
\vspace{0.15cm}
\item ¿Qué dimensión de la comunicación se fundamenta en el tono vocal, ritmo acústico y signos corporales no verbales?
\begin{enumerate}[label=\textbf{\Alph*)}]
\begin{multicols}{2}
  \item La lectura silenciosa
  \item La oralidad
  \item La tipografía digital
  \item La ortografía literal
\end{multicols}
\end{enumerate}
\vspace{0.15cm}
\item ¿Qué figura retórica magnifica o disminuye desproporcionadamente la realidad con fines expresivos o enfáticos?
\begin{enumerate}[label=\textbf{\Alph*)}]
\begin{multicols}{2}
  \item Hipérbole
  \item Asíndeton
  \item Elipsis
  \item Pleonasmo
\end{multicols}
\end{enumerate}
\vspace{0.15cm}
\item ¿Qué figura literaria consiste en atribuir propiedades, acciones o emociones humanas a seres inanimados o animales?
\begin{enumerate}[label=\textbf{\Alph*)}]
\begin{multicols}{2}
  \item Metonimia
  \item Personificación o prosopopeya
  \item Hipérbaton
  \item Sinestesia
\end{multicols}
\end{enumerate}
\vspace{0.15cm}
\end{enumerate}\n
\vspace{0.4cm}
\begin{tcolorbox}[colback=gray!5!white,colframe=gray!60!black,title=\textbf{CLAVE DE RESPUESTAS -- LITERATURA}]
\begin{multicols}{5}
\begin{enumerate}[label=\textbf{\arabic*.}]
\setcounter{enumi}{57}
  \item \textbf{B}
  \item \textbf{B}
  \item \textbf{B}
  \item \textbf{A}
  \item \textbf{B}
  \item \textbf{C}
  \item \textbf{B}
  \item \textbf{B}
  \item \textbf{B}
  \item \textbf{B}
  \item \textbf{B}
  \item \textbf{B}
  \item \textbf{B}
  \item \textbf{B}
  \item \textbf{B}
  \item \textbf{B}
  \item \textbf{B}
  \item \textbf{B}
  \item \textbf{B}
  \item \textbf{B}
  \item \textbf{A}
  \item \textbf{D}
  \item \textbf{B}
  \item \textbf{B}
  \item \textbf{B}
  \item \textbf{B}
  \item \textbf{B}
  \item \textbf{B}
  \item \textbf{B}
  \item \textbf{B}
  \item \textbf{B}
  \item \textbf{B}
  \item \textbf{B}
  \item \textbf{B}
  \item \textbf{B}
  \item \textbf{B}
  \item \textbf{B}
  \item \textbf{A}
  \item \textbf{B}
\end{enumerate}
\end{multicols}
\end{tcolorbox}
\newpage\n
% ==============================================================================
% MATERIA 4: GEOGRAFÍA DE EL SALVADOR Y CENTROAMÉRICA
% ==============================================================================
\section{4. Geografía de El Salvador y Centroamérica (50 Preguntas)}

\begin{enumerate}[label=\textbf{\arabic*.}]
\setcounter{enumi}{100}
\item ¿En cuántos departamentos se divide administrativamente el territorio de la República de El Salvador?
\begin{enumerate}[label=\textbf{\Alph*)}]
\begin{multicols}{2}
  \item 12 departamentos
  \item 14 departamentos
  \item 16 departamentos
  \item 18 departamentos
\end{multicols}
\end{enumerate}
\vspace{0.15cm}
\item ¿En qué departamento se sitúa el imponente volcán Chichontepec, popularmente apodado "Las Chiches" por su doble cumbre?
\begin{enumerate}[label=\textbf{\Alph*)}]
\begin{multicols}{2}
  \item Cabañas
  \item San Vicente
  \item Usulután
  \item La Paz
\end{multicols}
\end{enumerate}
\vspace{0.15cm}
\item ¿Cuál es el volcán de mayor altitud sobre el nivel del mar en la geografía salvadoreña?
\begin{enumerate}[label=\textbf{\Alph*)}]
\begin{multicols}{2}
  \item Volcán de San Salvador (Quezaltepeque)
  \item Volcán de Izalco
  \item Volcán de Santa Ana (Ilamatepec)
  \item Volcán de San Miguel (Chaparrastique)
\end{multicols}
\end{enumerate}
\vspace{0.15cm}
\item ¿En qué región geográfica de El Salvador se localiza el pintoresco circuito turístico de la "Ruta de las Flores"?
\begin{enumerate}[label=\textbf{\Alph*)}]
\begin{multicols}{2}
  \item Zona oriental
  \item Zona costera paracentral
  \item Zona occidental
  \item Zona norte fronteriza con Honduras
\end{multicols}
\end{enumerate}
\vspace{0.15cm}
\item ¿Qué departamentos salvadoreños atraviesa el circuito de pueblos coloniales de la Ruta de las Flores?
\begin{enumerate}[label=\textbf{\Alph*)}]
\begin{multicols}{2}
  \item San Salvador y La Libertad
  \item Ahuachapán y Sonsonate
  \item San Miguel y Morazán
  \item Cuscatlán y Cabañas
\end{multicols}
\end{enumerate}
\vspace{0.15cm}
\item ¿Qué sitio arqueológico salvadoreño destaca por sus templos prehispánicos en Chalchuapa, vecino a Joya de Cerén y San Andrés?
\begin{enumerate}[label=\textbf{\Alph*)}]
\begin{multicols}{2}
  \item Tazumal
  \item Quelepa
  \item Cihuatán
  \item Cara Sucia
\end{multicols}
\end{enumerate}
\vspace{0.15cm}
\item ¿Cuál es considerado el asentamiento arqueológico más representativo de la vida cotidiana prehispánica en El Salvador?
\begin{enumerate}[label=\textbf{\Alph*)}]
\begin{multicols}{2}
  \item Joya de Cerén
  \item El Trapiche
  \item Santa Leticia
  \item Corinto
\end{multicols}
\end{enumerate}
\vspace{0.15cm}
\item ¿Qué sitio arqueológico de El Salvador fue declarado Patrimonio Cultural de la Humanidad por la UNESCO en 1993?
\begin{enumerate}[label=\textbf{\Alph*)}]
\begin{multicols}{2}
  \item Tazumal
  \item Joya de Cerén
  \item San Andrés
  \item Tehuacán
\end{multicols}
\end{enumerate}
\vspace{0.15cm}
\item ¿Qué emblemático centro ceremonial maya ubicado en Chalchuapa cuenta con una gran pirámide escalonada de más de 20 metros de altura?
\begin{enumerate}[label=\textbf{\Alph*)}]
\begin{multicols}{2}
  \item Tazumal
  \item Madre Cacao
  \item Asunción
  \item El Cambio
\end{multicols}
\end{enumerate}
\vspace{0.15cm}
\item ¿Qué riqueza ecológica primordial resguarda el Parque Nacional El Imposible en el departamento de Ahuachapán?
\begin{enumerate}[label=\textbf{\Alph*)}]
\begin{multicols}{2}
  \item Arrecifes coralinos profundos
  \item La mayor biodiversidad de bosques y fauna silvestre amenazada del país
  \item Glaciares alpinos
  \item Desiertos de dunas
\end{multicols}
\end{enumerate}
\vspace{0.15cm}
\item ¿Qué parque nacional salvadoreño alberga un ecosistema único de bosque nebuloso en el punto de encuentro con Guatemala y Honduras?
\begin{enumerate}[label=\textbf{\Alph*)}]
\begin{multicols}{2}
  \item Parque Nacional Montecristo (Trifinio)
  \item Parque Walter Thilo Deininger
  \item Cerro Verde
  \item Conchagua
\end{multicols}
\end{enumerate}
\vspace{0.15cm}
\item ¿Qué parque natural es reconocido como uno de los últimos reductos del ecosistema de bosque seco y premontano en Mesoamérica?
\begin{enumerate}[label=\textbf{\Alph*)}]
\begin{multicols}{2}
  \item Parque Nacional El Imposible
  \item Laguna de Alegría
  \item El Boquerón
  \item La Puerta del Diablo
\end{multicols}
\end{enumerate}
\vspace{0.15cm}
\item ¿Qué reserva natural montañosa es ampliamente conocida por sus robles centenarios y bosques cubiertos de neblina permanente?
\begin{enumerate}[label=\textbf{\Alph*)}]
\begin{multicols}{2}
  \item Montecristo
  \item San Diego y San Felipe
  \item Nancuchiname
  \item Chaguantique
\end{multicols}
\end{enumerate}
\vspace{0.15cm}
\item ¿Cuál es la arteria fluvial más extensa, caudalosa e importante de la República de El Salvador?
\begin{enumerate}[label=\textbf{\Alph*)}]
\begin{multicols}{2}
  \item Río Paz
  \item Río Goascorán
  \item Río Lempa
  \item Río Grande de San Miguel
\end{multicols}
\end{enumerate}
\vspace{0.15cm}
\item ¿En qué departamento de la zona occidental se encuentra ubicado el paradisíaco Lago de Coatepeque?
\begin{enumerate}[label=\textbf{\Alph*)}]
\begin{multicols}{2}
  \item Santa Ana
  \item Sonsonate
  \item Ahuachapán
  \item La Libertad
\end{multicols}
\end{enumerate}
\vspace{0.15cm}
\item ¿Cuántos municipios integraron tradicionalmente la división territorial de El Salvador antes de la reforma de reestructuración municipal?
\begin{enumerate}[label=\textbf{\Alph*)}]
\begin{multicols}{2}
  \item 150 municipios
  \item 200 municipios
  \item 262 municipios (actualmente convertidos en distritos)
  \item 310 municipios
\end{multicols}
\end{enumerate}
\vspace{0.15cm}
\item ¿Qué departamento del occidente tiene por cabecera departamental a la "Ciudad Morena" homónima?
\begin{enumerate}[label=\textbf{\Alph*)}]
\begin{multicols}{2}
  \item Sonsonate
  \item Santa Ana
  \item Ahuachapán
  \item Chalatenango
\end{multicols}
\end{enumerate}
\vspace{0.15cm}
\item ¿Qué departamento montañoso del norte del país recibe el popular epíteto de "el departamento de las mil curvas"?
\begin{enumerate}[label=\textbf{\Alph*)}]
\begin{multicols}{2}
  \item Morazán
  \item Chalatenango
  \item Cabañas
  \item San Vicente
\end{multicols}
\end{enumerate}
\vspace{0.15cm}
\item ¿En qué departamento se ubica el Puerto de Acajutla, el principal puerto marítimo comercial del país sobre el Océano Pacífico?
\begin{enumerate}[label=\textbf{\Alph*)}]
\begin{multicols}{2}
  \item La Libertad
  \item Sonsonate
  \item La Paz
  \item La Unión
\end{multicols}
\end{enumerate}
\vspace{0.15cm}
\item ¿Cuál es el departamento con menor extensión territorial de toda la República de El Salvador?
\begin{enumerate}[label=\textbf{\Alph*)}]
\begin{multicols}{2}
  \item San Salvador
  \item Cuscatlán (aprox. 756 km²)
  \item Cabañas
  \item Ahuachapán
\end{multicols}
\end{enumerate}
\vspace{0.15cm}
\item ¿Qué prominente cráter volcánico y parque nacional se erige como emblema natural en las alturas de San Salvador?
\begin{enumerate}[label=\textbf{\Alph*)}]
\begin{multicols}{2}
  \item Cerro Guazapa
  \item Cerro El Boquerón (volcán de San Salvador)
  \item Cerro Cacahuatique
  \item Cerro Las Pavas
\end{multicols}
\end{enumerate}
\vspace{0.15cm}
\item ¿Qué playa del litoral de La Libertad ha adquirido renombre mundial como sede de campeonatos internacionales de surf?
\begin{enumerate}[label=\textbf{\Alph*)}]
\begin{multicols}{2}
  \item Playa El Espino
  \item Playa El Tunco (Surf City)
  \item Playa Metalío
  \item Playa Barra de Santiago
\end{multicols}
\end{enumerate}
\vspace{0.15cm}
\item ¿En qué departamento fronterizo se ubica el Cerro El Pital, la cumbre más alta de El Salvador a 2,730 msnm?
\begin{enumerate}[label=\textbf{\Alph*)}]
\begin{multicols}{2}
  \item Santa Ana
  \item Chalatenango
  \item Morazán
  \item San Miguel
\end{multicols}
\end{enumerate}
\vspace{0.15cm}
\item ¿Qué estratégica bahía y golfo del Océano Pacífico comparten tripartitamente El Salvador, Honduras y Nicaragua?
\begin{enumerate}[label=\textbf{\Alph*)}]
\begin{multicols}{2}
  \item Golfo de Nicoya
  \item Golfo de Fonseca
  \item Bahía de Jiquilisco
  \item Golfo de Tehuantepec
\end{multicols}
\end{enumerate}
\vspace{0.15cm}
\item ¿Cuál de las siguientes islas mayores del Golfo de Fonseca pertenece a la soberanía de El Salvador?
\begin{enumerate}[label=\textbf{\Alph*)}]
\begin{multicols}{2}
  \item Isla Tigre
  \item Isla Meanguera del Golfo
  \item Isla Zacate Grande
  \item Isla Ometepe
\end{multicols}
\end{enumerate}
\vspace{0.15cm}
\item ¿Qué parque arqueológico en el valle de Zapotitán es reconocido por su gran acrópolis y centro gubernamental maya?
\begin{enumerate}[label=\textbf{\Alph*)}]
\begin{multicols}{2}
  \item San Andrés
  \item Cara Sucia
  \item Quelepa
  \item Atiquizaya
\end{multicols}
\end{enumerate}
\vspace{0.15cm}
\item ¿En qué departamento oriental se ubica el trascendental sitio arqueológico lenca de Quelepa?
\begin{enumerate}[label=\textbf{\Alph*)}]
\begin{multicols}{2}
  \item Usulután
  \item San Miguel
  \item Morazán
  \item La Unión
\end{multicols}
\end{enumerate}
\vspace{0.15cm}
\item ¿Qué complejo marino-costero de Usulután es reconocido internacionalmente como sitio Ramsar y reserva de manglares?
\begin{enumerate}[label=\textbf{\Alph*)}]
\begin{multicols}{2}
  \item Bahía de Jiquilisco
  \item Estero de Jaltepeque
  \item Barra de Santiago
  \item Laguna del Jocotal
\end{multicols}
\end{enumerate}
\vspace{0.15cm}
\item ¿Qué volcán salvadoreño fue bautizado históricamente por los navegantes como "el Faro del Pacífico"?
\begin{enumerate}[label=\textbf{\Alph*)}]
\begin{multicols}{2}
  \item Volcán de San Vicente
  \item Volcán de Izalco
  \item Volcán Conchagua
  \item Volcán Tecapa
\end{multicols}
\end{enumerate}
\vspace{0.15cm}
\item ¿Qué majestuoso cuerpo lacustre se originó del colapso de dos conos volcánicos en Santa Ana?
\begin{enumerate}[label=\textbf{\Alph*)}]
\begin{multicols}{2}
  \item Lago de Ilopango
  \item Lago de Coatepeque
  \item Laguna de Güija
  \item Laguna de Olomega
\end{multicols}
\end{enumerate}
\vspace{0.15cm}
\item ¿Cuál es el cuerpo de agua artificial (embalse) de mayor extensión superficial en El Salvador?
\begin{enumerate}[label=\textbf{\Alph*)}]
\begin{multicols}{2}
  \item Lago de Ilopango
  \item Lago Suchitlán (Embalse Cerrón Grande, aprox. 135 km²)
  \item Embalse 15 de Septiembre
  \item Lago de Coatepeque
\end{multicols}
\end{enumerate}
\vspace{0.15cm}
\item ¿En qué departamento se concentra el Complejo Los Volcanes integrado por el Ilamatepec, Izalco y Cerro Verde?
\begin{enumerate}[label=\textbf{\Alph*)}]
\begin{multicols}{2}
  \item Santa Ana (y límites con Sonsonate)
  \item La Paz
  \item Cabañas
  \item Morazán
\end{multicols}
\end{enumerate}
\vspace{0.15cm}
\item ¿Qué destacada formación orográfica volcánica atraviesa el territorio salvadoreño de occidente al centro?
\begin{enumerate}[label=\textbf{\Alph*)}]
\begin{multicols}{2}
  \item Cordillera Apaneca-Ilamatepec
  \item Sierra de los Cuchumatanes
  \item Cordillera de Talamanca
  \item Sierra del Merendón
\end{multicols}
\end{enumerate}
\vspace{0.15cm}
\item ¿Qué importante cabecera de La Libertad es conocida cariñosamente como "la ciudad de las colinas"?
\begin{enumerate}[label=\textbf{\Alph*)}]
\begin{multicols}{2}
  \item Antiguo Cuscatlán
  \item Santa Tecla (Nueva San Salvador)
  \item Zaragoza
  \item San Juan Opico
\end{multicols}
\end{enumerate}
\vspace{0.15cm}
\item ¿Qué urbe es considerada el polo económico y demográfico más dinámico de la zona oriental salvadoreña?
\begin{enumerate}[label=\textbf{\Alph*)}]
\begin{multicols}{2}
  \item Usulután
  \item San Miguel
  \item La Unión
  \item San Francisco Gotera
\end{multicols}
\end{enumerate}
\vspace{0.15cm}
\item ¿En qué departamento se ubica el Aeropuerto Internacional de El Salvador San Óscar Arnulfo Romero y Galdámez?
\begin{enumerate}[label=\textbf{\Alph*)}]
\begin{multicols}{2}
  \item San Salvador
  \item La Paz (San Luis Talpa)
  \item La Libertad
  \item Cuscatlán
\end{multicols}
\end{enumerate}
\vspace{0.15cm}
\item ¿Cuál fue la cultura indígena y lingüística predominante que forjó la identidad en la mayor parte del territorio salvadoreño precolombino?
\begin{enumerate}[label=\textbf{\Alph*)}]
\begin{multicols}{2}
  \item Quechua
  \item Nahuat-Pipil
  \item Azteca del norte
  \item Tairona
\end{multicols}
\end{enumerate}
\vspace{0.15cm}
\item ¿Qué tradicional danza folclórica salvadoreña de origen colonial recrea la confrontación histórica entre cristianos y musulmanes?
\begin{enumerate}[label=\textbf{\Alph*)}]
\begin{multicols}{2}
  \item Danza del Torito Pinto
  \item Danza de los Moros y Cristianos (o Historiantes)
  \item Danza del tigre y el venado
  \item El Barreño
\end{multicols}
\end{enumerate}
\vspace{0.15cm}
\item ¿Qué expresión artesanal de modelado en barro cocido es de fama internacional en la ciudad de Ilobasco, Cabañas?
\begin{enumerate}[label=\textbf{\Alph*)}]
\begin{multicols}{2}
  \item Telares de hamacas
  \item Miniaturas y muñecos de barro tapados
  \item Talla de madera en nogal
  \item Cestería de mimbre
\end{multicols}
\end{enumerate}
\vspace{0.15cm}
\item ¿Cuál es el platillo gastronómico típico por excelencia y símbolo nacional de la gastronomía salvadoreña?
\begin{enumerate}[label=\textbf{\Alph*)}]
\begin{multicols}{2}
  \item Las pupusas
  \item Los tamales de elote
  \item Las enchiladas
  \item La yuca frita con chicharrón
\end{multicols}
\end{enumerate}
\vspace{0.15cm}
\item ¿Cuáles son las cuatro variedades de relleno más tradicionales e históricas de las pupusas salvadoreñas?
\begin{enumerate}[label=\textbf{\Alph*)}]
\begin{multicols}{2}
  \item Queso, frijol con queso, revueltas (chicharrón) y ayote
  \item Jamón, tocino, pepperoni y piña
  \item Pescado, mariscos, cangrejo y pulpo
  \item Carne asada, pollo, chorizo y res
\end{multicols}
\end{enumerate}
\vspace{0.15cm}
\item ¿Qué ancestral bebida fermentada de maíz, dulce de panela y jengibre se elabora en fiestas tradicionales salvadoreñas?
\begin{enumerate}[label=\textbf{\Alph*)}]
\begin{multicols}{2}
  \item Chicha
  \item Horchata de morro
  \item Cebada
  \item Fresco de ensalada
\end{multicols}
\end{enumerate}
\vspace{0.15cm}
\item ¿Qué fruto agridulce con hueso leñoso, del género Spondias, es autóctono y sumamente popular en El Salvador (en miel o verde)?
\begin{enumerate}[label=\textbf{\Alph*)}]
\begin{multicols}{2}
  \item Manzana
  \item Jocote (corona, de invierno, etc.)
  \item Uva
  \item Pera
\end{multicols}
\end{enumerate}
\vspace{0.15cm}
\item ¿En qué mes del año se celebran las Fiestas Patronales de San Salvador en honor al Divino Salvador del Mundo?
\begin{enumerate}[label=\textbf{\Alph*)}]
\begin{multicols}{2}
  \item Mayo
  \item Agosto (del 1 al 6 de agosto)
  \item Septiembre
  \item Diciembre
\end{multicols}
\end{enumerate}
\vspace{0.15cm}
\item ¿Qué manifestación de devoción popular y arte efímero engalana las calles salvadoreñas durante el Viernes Santo?
\begin{enumerate}[label=\textbf{\Alph*)}]
\begin{multicols}{2}
  \item Alfombras artesanales de aserrín teñido y sal para las procesiones del Santo Entierro
  \item Desfiles militares
  \item Carnavales de carrozas
  \item Feria del café
\end{multicols}
\end{enumerate}
\vspace{0.15cm}
\item ¿Qué instrumento de cuerdas pulsadas es indispensable en los tríos y conjuntos de música folclórica y campesina salvadoreña?
\begin{enumerate}[label=\textbf{\Alph*)}]
\begin{multicols}{2}
  \item Gaita
  \item Guitarra acústica
  \item Sitar
  \item Balalaika
\end{multicols}
\end{enumerate}
\vspace{0.15cm}
\item ¿En qué ciudad occidental natal de Alfredo Espino se conserva la Casa Museo que honra su memoria?
\begin{enumerate}[label=\textbf{\Alph*)}]
\begin{multicols}{2}
  \item Ahuachapán
  \item Metapán
  \item Izalco
  \item Chalchuapa
\end{multicols}
\end{enumerate}
\vspace{0.15cm}
\item ¿Qué prominente formación montañosa al norte de San Salvador fue un emblemático bastión insurgente durante la guerra civil?
\begin{enumerate}[label=\textbf{\Alph*)}]
\begin{multicols}{2}
  \item Cerro Verde
  \item Cerro de Guazapa
  \item Cerro El Boquerón
  \item Cerro de las Chiches
\end{multicols}
\end{enumerate}
\vspace{0.15cm}
\item ¿Cerca de qué municipio al norte de San Salvador se localiza el parque arqueológico de la gran ciudadela de Cihuatán?
\begin{enumerate}[label=\textbf{\Alph*)}]
\begin{multicols}{2}
  \item Aguilares
  \item Apopa
  \item San Martín
  \item Nejapa
\end{multicols}
\end{enumerate}
\vspace{0.15cm}
\item ¿Qué importante territorio natural salvadoreño fue designado por la UNESCO como Reserva de la Biosfera en 2007?
\begin{enumerate}[label=\textbf{\Alph*)}]
\begin{multicols}{2}
  \item Reserva de la Biosfera Apaneca-Ilamatepec
  \item Parque Cuscatlán
  \item Bosque San Andrés
  \item Laguna de Ilopango
\end{multicols}
\end{enumerate}
\vspace{0.15cm}
\end{enumerate}\n
\vspace{0.4cm}
\begin{tcolorbox}[colback=gray!5!white,colframe=gray!60!black,title=\textbf{CLAVE DE RESPUESTAS -- GEOGRAFÍA DE EL SALVADOR Y CENTROAMÉRICA}]
\begin{multicols}{5}
\begin{enumerate}[label=\textbf{\arabic*.}]
\setcounter{enumi}{100}
  \item \textbf{B}
  \item \textbf{B}
  \item \textbf{C}
  \item \textbf{C}
  \item \textbf{B}
  \item \textbf{A}
  \item \textbf{A}
  \item \textbf{B}
  \item \textbf{A}
  \item \textbf{B}
  \item \textbf{A}
  \item \textbf{A}
  \item \textbf{A}
  \item \textbf{C}
  \item \textbf{A}
  \item \textbf{C}
  \item \textbf{B}
  \item \textbf{B}
  \item \textbf{B}
  \item \textbf{B}
  \item \textbf{B}
  \item \textbf{B}
  \item \textbf{B}
  \item \textbf{B}
  \item \textbf{B}
  \item \textbf{A}
  \item \textbf{B}
  \item \textbf{A}
  \item \textbf{B}
  \item \textbf{B}
  \item \textbf{B}
  \item \textbf{A}
  \item \textbf{A}
  \item \textbf{B}
  \item \textbf{B}
  \item \textbf{B}
  \item \textbf{B}
  \item \textbf{B}
  \item \textbf{B}
  \item \textbf{A}
  \item \textbf{A}
  \item \textbf{A}
  \item \textbf{B}
  \item \textbf{B}
  \item \textbf{A}
  \item \textbf{B}
  \item \textbf{A}
  \item \textbf{B}
  \item \textbf{A}
  \item \textbf{A}
\end{enumerate}
\end{multicols}
\end{tcolorbox}
\newpage\n
% ==============================================================================
% MATERIA 5: REALIDAD ECONÓMICA Y SOCIAL
% ==============================================================================
\section{5. Realidad Económica y Social (9 Preguntas)}

\begin{enumerate}[label=\textbf{\arabic*.}]
\setcounter{enumi}{150}
\item ¿Cuál es la tasa porcentual general vigente que se aplica al Impuesto al Valor Agregado (IVA) en El Salvador?
\begin{enumerate}[label=\textbf{\Alph*)}]
\begin{multicols}{2}
  \item 10\%
  \item 12\%
  \item 13\%
  \item 15\%
\end{multicols}
\end{enumerate}
\vspace{0.15cm}
\item ¿Qué tipo de gravamen indirecto representa la tasa del 13\% en la adquisición de bienes de consumo en El Salvador?
\begin{enumerate}[label=\textbf{\Alph*)}]
\begin{multicols}{2}
  \item Tasa fija general del IVA
  \item Impuesto sobre la Renta
  \item Arancel de aduanas
  \item Impuesto territorial
\end{multicols}
\end{enumerate}
\vspace{0.15cm}
\item ¿Qué operaciones comerciales y productos cuentan con tasa preferencial del 0\% o exención expresa de IVA en El Salvador?
\begin{enumerate}[label=\textbf{\Alph*)}]
\begin{multicols}{2}
  \item Exportaciones, medicinas esenciales y productos de canasta básica exentos
  \item Vehículos de lujo
  \item Bebidas alcohólicas importadas
  \item Telefonía celular corporativa
\end{multicols}
\end{enumerate}
\vspace{0.15cm}
\item Desde el punto de vista socioeconómico, ¿con qué carencia primordial se relaciona la condición de pobreza monetaria de un hogar?
\begin{enumerate}[label=\textbf{\Alph*)}]
\begin{multicols}{2}
  \item Falta de ingresos económicos suficientes para cubrir la Canasta Básica Alimentaria y de servicios
  \item Ausencia de vehículos propios
  \item Falta de entretenimiento digital
  \item Vivir lejos de la capital
\end{multicols}
\end{enumerate}
\vspace{0.15cm}
\item Según informes económicos de la CEPAL (2023), ¿qué porcentaje aproximado de la población salvadoreña se situaba bajo la línea de pobreza?
\begin{enumerate}[label=\textbf{\Alph*)}]
\begin{multicols}{2}
  \item 12.4\%
  \item 26.6\%
  \item 45.0\%
  \item 60.2\%
\end{multicols}
\end{enumerate}
\vspace{0.15cm}
\item De acuerdo con el Banco Mundial, ¿qué tasa aproximada reflejó el índice de pobreza extrema en El Salvador para el año 2023?
\begin{enumerate}[label=\textbf{\Alph*)}]
\begin{multicols}{2}
  \item Alrededor del 9.3\%
  \item Menos del 1\%
  \item Más del 35\%
  \item Exactamente el 20\%
\end{multicols}
\end{enumerate}
\vspace{0.15cm}
\item ¿Cuál es la función primordial de una Administradora de Fondos de Pensiones (AFP)?
\begin{enumerate}[label=\textbf{\Alph*)}]
\begin{multicols}{2}
  \item Conceder créditos automotrices
  \item Administrar e invertir las cuentas de ahorro individual de los trabajadores para su futura jubilación
  \item Cobrar los impuestos municipales
  \item Pagar las deudas del Estado
\end{multicols}
\end{enumerate}
\vspace{0.15cm}
\item ¿Qué significado tienen las siglas financieras AFP en el sistema de seguridad social salvadoreño?
\begin{enumerate}[label=\textbf{\Alph*)}]
\begin{multicols}{2}
  \item Asociación Federal de Préstamos
  \item Administradora de Fondos de Pensiones
  \item Agencia Fiscal Previsional
  \item Asamblea Financiera Pública
\end{multicols}
\end{enumerate}
\vspace{0.15cm}
\item ¿Qué porcentaje de su salario mensual aporta obligatoriamente el trabajador al Sistema de Ahorro para Pensiones (AFP)?
\begin{enumerate}[label=\textbf{\Alph*)}]
\begin{multicols}{2}
  \item 3.00\%
  \item 5.50\%
  \item 7.25\%
  \item 10.00\%
\end{multicols}
\end{enumerate}
\vspace{0.15cm}
\end{enumerate}\n
\vspace{0.4cm}
\begin{tcolorbox}[colback=gray!5!white,colframe=gray!60!black,title=\textbf{CLAVE DE RESPUESTAS -- REALIDAD ECONÓMICA Y SOCIAL}]
\begin{multicols}{5}
\begin{enumerate}[label=\textbf{\arabic*.}]
\setcounter{enumi}{150}
  \item \textbf{C}
  \item \textbf{A}
  \item \textbf{A}
  \item \textbf{A}
  \item \textbf{B}
  \item \textbf{A}
  \item \textbf{B}
  \item \textbf{B}
  \item \textbf{C}
\end{enumerate}
\end{multicols}
\end{tcolorbox}
\newpage\n
% ==============================================================================
% MATERIA 6: CONSTITUCIÓN DE LA REPÚBLICA Y DDHH
% ==============================================================================
\section{6. Constitución de la República y DDHH (27 Preguntas)}

\begin{enumerate}[label=\textbf{\arabic*.}]
\setcounter{enumi}{159}
\item El artículo 85 de la Constitución consagra que El Salvador es una República cuya forma de gobierno es:
\begin{enumerate}[label=\textbf{\Alph*)}]
\begin{multicols}{2}
  \item Monárquica parlamentaria
  \item Republicana, democrática y representativa
  \item Teocrática
  \item Militar dictatorial
\end{multicols}
\end{enumerate}
\vspace{0.15cm}
\item Según los principios constitucionales fundamentales de la República de El Salvador, el Estado es:
\begin{enumerate}[label=\textbf{\Alph*)}]
\begin{multicols}{2}
  \item Soberano, democrático e independiente
  \item Sujeto a mandatos extranjeros
  \item Una confederación de clanes
  \item Un protectorado internacional
\end{multicols}
\end{enumerate}
\vspace{0.15cm}
\item ¿Qué régimen político rige en El Salvador en cuanto a la concentración de la jefatura de Estado y de Gobierno en el Presidente de la República?
\begin{enumerate}[label=\textbf{\Alph*)}]
\begin{multicols}{2}
  \item Sistema parlamentario
  \item Sistema presidencialista
  \item Directorio colegiado
  \item Monarquía constitucional
\end{multicols}
\end{enumerate}
\vspace{0.15cm}
\item ¿En qué cuerpo normativo fundamental del ordenamiento jurídico se consagran los derechos individuales de los salvadoreños?
\begin{enumerate}[label=\textbf{\Alph*)}]
\begin{multicols}{2}
  \item En el Código de Comercio
  \item En la Constitución de la República (Título II)
  \item En la Ley de Imprenta
  \item En las Ordenanzas Municipales
\end{multicols}
\end{enumerate}
\vspace{0.15cm}
\item ¿Qué artículo de la Constitución de El Salvador garantiza el derecho de toda persona a expresar y difundir libremente sus pensamientos?
\begin{enumerate}[label=\textbf{\Alph*)}]
\begin{multicols}{2}
  \item Artículo 2
  \item Artículo 6
  \item Artículo 12
  \item Artículo 25
\end{multicols}
\end{enumerate}
\vspace{0.15cm}
\item ¿Qué garantía y proceso constitucional permite a una persona protegerse contra privaciones ilegales o arbitrarias de su libertad física?
\begin{enumerate}[label=\textbf{\Alph*)}]
\begin{multicols}{2}
  \item El Hábeas Corpus (o Exhibición Personal)
  \item El Recurso de Apelación
  \item El Juicio de Cuentas
  \item La Casación
\end{multicols}
\end{enumerate}
\vspace{0.15cm}
\item ¿Cuál de los siguientes términos NO constituye un derecho humano individual reconocido en el orden constitucional salvadoreño?
\begin{enumerate}[label=\textbf{\Alph*)}]
\begin{multicols}{2}
  \item Derecho a la vida y a la integridad física
  \item Derecho canónico (normativa eclesiástica interna de una iglesia)
  \item Derecho a la seguridad y a la propiedad
  \item Derecho a la libertad de circulación
\end{multicols}
\end{enumerate}
\vspace{0.15cm}
\item En el ordenamiento constitucional de El Salvador, el acto de votar en elecciones populares constituye:
\begin{enumerate}[label=\textbf{\Alph*)}]
\begin{multicols}{2}
  \item Una actividad comercial
  \item Un derecho y un deber político ciudadano
  \item Una obligación impuesta solo a funcionarios
  \item Un trámite optativo para extranjeros
\end{multicols}
\end{enumerate}
\vspace{0.15cm}
\item Según el artículo 78 de la Constitución, ¿qué características inviolables definen al sufragio en El Salvador?
\begin{enumerate}[label=\textbf{\Alph*)}]
\begin{multicols}{2}
  \item Censitario, público y delegado
  \item Universal, directo, igualitario y secreto
  \item Condicionado por la renta
  \item Exclusivo para varones
\end{multicols}
\end{enumerate}
\vspace{0.15cm}
\item ¿Qué artículo de la Constitución salvadoreña enuncia los derechos políticos fundamentales de los ciudadanos, entre ellos ejercer el sufragio?
\begin{enumerate}[label=\textbf{\Alph*)}]
\begin{multicols}{2}
  \item Artículo 72
  \item Artículo 14
  \item Artículo 38
  \item Artículo 101
\end{multicols}
\end{enumerate}
\vspace{0.15cm}
\item ¿Qué principio democrático medular garantiza que la decisión electoral del ciudadano se emita sin coacción ni conocimiento ajeno?
\begin{enumerate}[label=\textbf{\Alph*)}]
\begin{multicols}{2}
  \item Principio de notoriedad
  \item Principio de libertad y secreto del voto
  \item Principio de jerarquía militar
  \item Principio de delegación obligada
\end{multicols}
\end{enumerate}
\vspace{0.15cm}
\item ¿Cuál es una atribución constitucional central encomendada al Ministerio Público salvadoreño (FGR, PDDH, PGR)?
\begin{enumerate}[label=\textbf{\Alph*)}]
\begin{multicols}{2}
  \item Legislar leyes tributarias
  \item Investigar los delitos y promover la acción de la justicia en defensa de la legalidad
  \item Dictar sentencias penales definitivas
  \item Comandar la Fuerza Armada
\end{multicols}
\end{enumerate}
\vspace{0.15cm}
\item ¿Qué labor esencial distingue a la Fiscalía General de la República dentro del proceso penal salvadoreño?
\begin{enumerate}[label=\textbf{\Alph*)}]
\begin{multicols}{2}
  \item Administrar las cárceles
  \item Dirigir la investigación del delito y ejercer la acción penal pública (acusación)
  \item Juzgar y condenar a los procesados
  \item Custodiar las fronteras aéreas
\end{multicols}
\end{enumerate}
\vspace{0.15cm}
\item ¿Qué funcionario electo por la Asamblea Legislativa ostenta la máxima autoridad y dirección de la Fiscalía General de la República?
\begin{enumerate}[label=\textbf{\Alph*)}]
\begin{multicols}{2}
  \item Presidente de la Corte Suprema
  \item Fiscal General de la República
  \item Procurador para la Defensa de los DDHH
  \item Ministro de Justicia
\end{multicols}
\end{enumerate}
\vspace{0.15cm}
\item ¿Cuál es la misión institucional principal de la Policía Nacional Civil (PNC) en la sociedad salvadoreña?
\begin{enumerate}[label=\textbf{\Alph*)}]
\begin{multicols}{2}
  \item Prevenir, combatir e investigar los hechos delictivos y garantizar el orden y la seguridad ciudadana
  \item Dictar leyes sobre comercio
  \item Recaudar impuestos de importación
  \item Representar al país diplomáticamente
\end{multicols}
\end{enumerate}
\vspace{0.15cm}
\item ¿En qué año se fundó la Policía Nacional Civil (PNC) como institución civil desmilitarizada emanada de los Acuerdos de Paz?
\begin{enumerate}[label=\textbf{\Alph*)}]
\begin{multicols}{2}
  \item 1980
  \item 1992
  \item 1998
  \item 2001
\end{multicols}
\end{enumerate}
\vspace{0.15cm}
\item ¿Cuál es el lema institucional que enarbola la Policía Nacional Civil de El Salvador?
\begin{enumerate}[label=\textbf{\Alph*)}]
\begin{multicols}{2}
  \item "Honor y Gloria"
  \item "Servicio - Orden - Seguridad"
  \item "Dios, Unión y Libertad"
  \item "Lealtad y Sacrificio"
\end{multicols}
\end{enumerate}
\vspace{0.15cm}
\item ¿Qué rango jerárquico es el más alto en la escala de mando y carrera de la Policía Nacional Civil?
\begin{enumerate}[label=\textbf{\Alph*)}]
\begin{multicols}{2}
  \item Subcomisionado
  \item Comisionado General
  \item Inspector Jefe
  \item Sargento Mayor
\end{multicols}
\end{enumerate}
\vspace{0.15cm}
\item ¿Cuál es la función primordial y carácter esencial asignado a la Fuerza Armada por la Constitución salvadoreña?
\begin{enumerate}[label=\textbf{\Alph*)}]
\begin{multicols}{2}
  \item Defender la soberanía del Estado y la integridad de su territorio
  \item Manejar los bancos comerciales
  \item Intervenir en disputas de partidos políticos
  \item Nombrar a los diputados
\end{multicols}
\end{enumerate}
\vspace{0.15cm}
\item El artículo 212 de la Constitución estipula que la misión militar de la Fuerza Armada salvadoreña es:
\begin{enumerate}[label=\textbf{\Alph*)}]
\begin{multicols}{2}
  \item Garantizar la alternabilidad en el poder municipal
  \item Defender la soberanía del Estado y la integridad del territorio
  \item Ejercer el comercio exterior
  \item Investigar delitos comunes
\end{multicols}
\end{enumerate}
\vspace{0.15cm}
\item ¿Cuál de los siguientes cuerpos NO es una de las ramas orgánicas de la Fuerza Armada de El Salvador?
\begin{enumerate}[label=\textbf{\Alph*)}]
\begin{multicols}{2}
  \item Fuerza de Seguridad Pública (antigua denominación policial disuelta)
  \item Ejército
  \item Fuerza Aérea
  \item Fuerza Naval
\end{multicols}
\end{enumerate}
\vspace{0.15cm}
\item En el marco del derecho penal y administrativo salvadoreño, ¿quién tiene legalmente la investidura de "agente de autoridad"?
\begin{enumerate}[label=\textbf{\Alph*)}]
\begin{multicols}{2}
  \item Un comerciante particular
  \item Un agente de la Policía Nacional Civil
  \item Un vigilante privado sin acreditación
  \item Un estudiante de secundaria
\end{multicols}
\end{enumerate}
\vspace{0.15cm}
\item ¿Bajo qué circunstancias excepcionales puede un agente de la PNC allanar una morada o domicilio sin previa orden judicial escrita?
\begin{enumerate}[label=\textbf{\Alph*)}]
\begin{multicols}{2}
  \item En cualquier momento por mera sospecha
  \item Únicamente en caso de flagrancia delictiva, persecución urgente o grave peligro de vidas (emergencia)
  \item Para cobrar un préstamo vencido
  \item Por solicitud de un vecino
\end{multicols}
\end{enumerate}
\vspace{0.15cm}
\item ¿Qué cuerpo normativo especial regula las obligaciones, derechos, ascensos y régimen disciplinario de los agentes de la PNC?
\begin{enumerate}[label=\textbf{\Alph*)}]
\begin{multicols}{2}
  \item Código Aduanero
  \item Ley de la Carrera Policial
  \item Ley de Notariado
  \item Ley de Inquilinato
\end{multicols}
\end{enumerate}
\vspace{0.15cm}
\item ¿Cuál es el fin universal y supremo que persigue el reconocimiento de los Derechos Humanos?
\begin{enumerate}[label=\textbf{\Alph*)}]
\begin{multicols}{2}
  \item Proteger la dignidad intrínseca de todas las personas frente a abusos de poder
  \item Limitar el crecimiento poblacional
  \item Establecer tasas comerciales entre países
  \item Facilitar la adquisición de bienes
\end{multicols}
\end{enumerate}
\vspace{0.15cm}
\item ¿Qué histórico documento internacional adoptado por la ONU en 1948 es la piedra angular de los Derechos Humanos modernos?
\begin{enumerate}[label=\textbf{\Alph*)}]
\begin{multicols}{2}
  \item Pacto de Varsovia
  \item Declaración Universal de los Derechos Humanos
  \item Carta de las Américas
  \item Tratado de Versalles
\end{multicols}
\end{enumerate}
\vspace{0.15cm}
\item ¿Qué texto constitucional salvadoreño del siglo XIX fue pionero en consagrar un catálogo amplio de libertades individuales y laicismo?
\begin{enumerate}[label=\textbf{\Alph*)}]
\begin{multicols}{2}
  \item Constitución de 1824
  \item Constitución de 1886
  \item Constitución de 1939
  \item Constitución de 1950
\end{multicols}
\end{enumerate}
\vspace{0.15cm}
\end{enumerate}\n
\vspace{0.4cm}
\begin{tcolorbox}[colback=gray!5!white,colframe=gray!60!black,title=\textbf{CLAVE DE RESPUESTAS -- CONSTITUCIÓN DE LA REPÚBLICA Y DDHH}]
\begin{multicols}{5}
\begin{enumerate}[label=\textbf{\arabic*.}]
\setcounter{enumi}{159}
  \item \textbf{B}
  \item \textbf{A}
  \item \textbf{B}
  \item \textbf{B}
  \item \textbf{B}
  \item \textbf{A}
  \item \textbf{B}
  \item \textbf{B}
  \item \textbf{B}
  \item \textbf{A}
  \item \textbf{B}
  \item \textbf{B}
  \item \textbf{B}
  \item \textbf{B}
  \item \textbf{A}
  \item \textbf{B}
  \item \textbf{B}
  \item \textbf{B}
  \item \textbf{A}
  \item \textbf{B}
  \item \textbf{A}
  \item \textbf{B}
  \item \textbf{B}
  \item \textbf{B}
  \item \textbf{A}
  \item \textbf{B}
  \item \textbf{B}
\end{enumerate}
\end{multicols}
\end{tcolorbox}
\newpage\n
% ==============================================================================
% MATERIA 7: LEGISLACIÓN SALVADOREÑA
% ==============================================================================
\section{7. Legislación Salvadoreña (19 Preguntas)}

\begin{enumerate}[label=\textbf{\arabic*.}]
\setcounter{enumi}{186}
\item ¿Qué es el Código Penal dentro del sistema de leyes salvadoreño?
\begin{enumerate}[label=\textbf{\Alph*)}]
\begin{multicols}{2}
  \item Una ley secundaria que tipifica las conductas delictivas, las faltas y las penas correspondientes
  \item Un manual de tránsito vial
  \item Un reglamento municipal de mercados
  \item Un tratado de comercio exterior
\end{multicols}
\end{enumerate}
\vspace{0.15cm}
\item ¿Cuál de las siguientes figuras típicas está catalogada como delito grave contra el bien jurídico más preciado: la vida?
\begin{enumerate}[label=\textbf{\Alph*)}]
\begin{multicols}{2}
  \item Perturbación del orden público menor
  \item Homicidio (en sus distintas modalidades)
  \item Estacionar en zona prohibida
  \item Girar un cheque con insuficiencia temporal
\end{multicols}
\end{enumerate}
\vspace{0.15cm}
\item ¿Cuál es la pena máxima de prisión que contempla el Código Penal salvadoreño para el delito de Homicidio Agravado?
\begin{enumerate}[label=\textbf{\Alph*)}]
\begin{multicols}{2}
  \item 20 años
  \item 30 años
  \item Hasta 50 años de prisión
  \item Pena perpetua fija
\end{multicols}
\end{enumerate}
\vspace{0.15cm}
\item ¿Qué función primordial cumple el Código Procesal Penal de la República de El Salvador?
\begin{enumerate}[label=\textbf{\Alph*)}]
\begin{multicols}{2}
  \item Establecer las normas de tránsito vehicular
  \item Regular el conjunto de normas y trámites que rigen el proceso penal para investigar y juzgar los delitos
  \item Cobrar los impuestos a la importación
  \item Fijar los feriados nacionales
\end{multicols}
\end{enumerate}
\vspace{0.15cm}
\item ¿Qué fase definitoria y pública regula el Código Procesal Penal para la valoración de pruebas por los jueces?
\begin{enumerate}[label=\textbf{\Alph*)}]
\begin{multicols}{2}
  \item El arbitraje comercial
  \item El juicio plenario u oral y público
  \item La mediación laboral
  \item La asamblea de socios
\end{multicols}
\end{enumerate}
\vspace{0.15cm}
\item ¿En qué año entró en vigencia en El Salvador la reforma procesal que instauró el moderno sistema penal acusatorio oral?
\begin{enumerate}[label=\textbf{\Alph*)}]
\begin{multicols}{2}
  \item 1980
  \item 1992
  \item 1998
  \item 2005
\end{multicols}
\end{enumerate}
\vspace{0.15cm}
\item ¿Qué principio procesal y constitucional asegura que toda persona es considerada inocente mientras no se declare su culpabilidad en juicio?
\begin{enumerate}[label=\textbf{\Alph*)}]
\begin{multicols}{2}
  \item Principio de culpabilidad presunta
  \item Principio de presunción de inocencia (Art. 12 Cn)
  \item Principio de flagrancia
  \item Principio inquisitorio
\end{multicols}
\end{enumerate}
\vspace{0.15cm}
\item Según la Ley de Tránsito salvadoreña, ¿qué requisito de licencia es indispensable para conducir vehículos de transporte pesado de carga?
\begin{enumerate}[label=\textbf{\Alph*)}]
\begin{multicols}{2}
  \item Poseer licencia particular tipo C
  \item Poseer licencia de conducir tipo T (transporte pesado)
  \item Licencia de motocicleta juvenil
  \item Carnet de estudiante universitario
\end{multicols}
\end{enumerate}
\vspace{0.15cm}
\item ¿A qué consecuencia penal se expone un conductor que maneje un vehículo en evidente estado de ebriedad según el Código Penal?
\begin{enumerate}[label=\textbf{\Alph*)}]
\begin{multicols}{2}
  \item Solo una amonestación verbal
  \item Detención y pena de prisión por el delito de conducción peligrosa de vehículos de motor
  \item Pagar el combustible de la patrulla
  \item Ninguna sanción si conduce despacio
\end{multicols}
\end{enumerate}
\vspace{0.15cm}
\item ¿Qué conjunto de sanciones administrativas y judiciales recibe un conductor que es sorprendido manejando ebrio en El Salvador?
\begin{enumerate}[label=\textbf{\Alph*)}]
\begin{multicols}{2}
  \item Multa muy grave
  \item Decomiso del vehículo y retención de la licencia de conducir
  \item Remisión al juzgado de paz por delito penal
  \item Todas las anteriores
\end{multicols}
\end{enumerate}
\vspace{0.15cm}
\item ¿Cuál de los siguientes es un requisito indispensable exigido por la ANSP para ingresar como aspirante a la carrera policial?
\begin{enumerate}[label=\textbf{\Alph*)}]
\begin{multicols}{2}
  \item Tener experiencia militar previa
  \item Carecer de antecedentes penales y policiales y superar las pruebas de confiabilidad
  \item Ser propietario de una vivienda
  \item Estar afiliado a un partido político
\end{multicols}
\end{enumerate}
\vspace{0.15cm}
\item ¿Qué propósito primordial persigue la Ley de la Carrera Policial en el fortalecimiento institucional?
\begin{enumerate}[label=\textbf{\Alph*)}]
\begin{multicols}{2}
  \item Militarizar la administración pública
  \item Profesionalizar y tecnificar a los miembros de la Policía Nacional Civil
  \item Reducir el salario de los agentes
  \item Eliminar la escala de ascensos
\end{multicols}
\end{enumerate}
\vspace{0.15cm}
\item Según la Ley Orgánica de la PNC y la Constitución, ¿qué autoridad nombra al Director General de la corporación?
\begin{enumerate}[label=\textbf{\Alph*)}]
\begin{multicols}{2}
  \item La Asamblea Legislativa
  \item El Presidente de la República
  \item La Corte Suprema de Justicia
  \item El Tribunal Supremo Electoral
\end{multicols}
\end{enumerate}
\vspace{0.15cm}
\item ¿Cuál es el fin supremo que define la actuación operativa y moral de la Policía Nacional Civil?
\begin{enumerate}[label=\textbf{\Alph*)}]
\begin{multicols}{2}
  \item El cobro coactivo de impuestos
  \item El mantenimiento de la paz, la tranquilidad social y el orden público
  \item La construcción de carreteras
  \item La explotación de recursos minerales
\end{multicols}
\end{enumerate}
\vspace{0.15cm}
\item ¿Cuál es la finalidad principal de la moderna "Ley Crecer Juntos para la Protección Integral de la Primera Infancia, Niñez y Adolescencia"?
\begin{enumerate}[label=\textbf{\Alph*)}]
\begin{multicols}{2}
  \item Promover la contratación laboral de menores
  \item Garantizar el ejercicio y disfrute pleno de los derechos de toda niña, niño y adolescente
  \item Crear cárceles juveniles
  \item Regular el matrimonio de jóvenes
\end{multicols}
\end{enumerate}
\vspace{0.15cm}
\item ¿A qué sector vulnerable de la población salvadoreña tutela prioritariamente la Ley Crecer Juntos?
\begin{enumerate}[label=\textbf{\Alph*)}]
\begin{multicols}{2}
  \item Empresarios exportadores
  \item Primera infancia, niñas, niños y adolescentes
  \item Funcionarios públicos diplomáticos
  \item Propietarios de bienes raíces
\end{multicols}
\end{enumerate}
\vspace{0.15cm}
\item ¿Qué institución estatal especializada es el ente rector para aplicar y velar por la Ley Crecer Juntos?
\begin{enumerate}[label=\textbf{\Alph*)}]
\begin{multicols}{2}
  \item CONNA (actualmente transformado en CONAPINA)
  \item Ministerio de Hacienda
  \item Alcaldía de San Salvador
  \item Corte de Cuentas
\end{multicols}
\end{enumerate}
\vspace{0.15cm}
\item ¿En qué articulado de la Ley Orgánica de la PNC se definen los deberes y principios básicos de actuación policial?
\begin{enumerate}[label=\textbf{\Alph*)}]
\begin{multicols}{2}
  \item Artículo 2 (y principios deontológicos institucionales)
  \item Artículo 95
  \item Artículo 150
  \item Artículo 200
\end{multicols}
\end{enumerate}
\vspace{0.15cm}
\item ¿Qué ley especial regula el proceso técnico de selección, ingreso, formación en la ANSP y ascensos en la PNC?
\begin{enumerate}[label=\textbf{\Alph*)}]
\begin{multicols}{2}
  \item Código Tributario
  \item Ley de la Carrera Policial
  \item Ley de Contrataciones del Estado
  \item Código de Familia
\end{multicols}
\end{enumerate}
\vspace{0.15cm}
\end{enumerate}\n
\vspace{0.4cm}
\begin{tcolorbox}[colback=gray!5!white,colframe=gray!60!black,title=\textbf{CLAVE DE RESPUESTAS -- LEGISLACIÓN SALVADOREÑA}]
\begin{multicols}{5}
\begin{enumerate}[label=\textbf{\arabic*.}]
\setcounter{enumi}{186}
  \item \textbf{A}
  \item \textbf{B}
  \item \textbf{C}
  \item \textbf{B}
  \item \textbf{B}
  \item \textbf{C}
  \item \textbf{B}
  \item \textbf{B}
  \item \textbf{B}
  \item \textbf{D}
  \item \textbf{B}
  \item \textbf{B}
  \item \textbf{B}
  \item \textbf{B}
  \item \textbf{B}
  \item \textbf{B}
  \item \textbf{A}
  \item \textbf{A}
  \item \textbf{B}
\end{enumerate}
\end{multicols}
\end{tcolorbox}
\newpage\n
% ==============================================================================
% MATERIA 8: MATEMÁTICAS
% ==============================================================================
\section{8. Matemáticas (25 Preguntas)}

\begin{enumerate}[label=\textbf{\arabic*.}]
\setcounter{enumi}{205}
\item Resuelve la ecuación lineal de primer grado: \(x + 3x + 100 = 500\)
\begin{enumerate}[label=\textbf{\Alph*)}]
\begin{multicols}{2}
  \item \(x = 50\)
  \item \(x = 100\)
  \item \(x = 125\)
  \item \(x = 150\)
\end{multicols}
\end{enumerate}
\vspace{0.15cm}
\item Encuentra el valor de \(x\) que satisface la igualdad: \(2x + 5 = 15\)
\begin{enumerate}[label=\textbf{\Alph*)}]
\begin{multicols}{2}
  \item \(x = 3\)
  \item \(x = 5\)
  \item \(x = 7\)
  \item \(x = 10\)
\end{multicols}
\end{enumerate}
\vspace{0.15cm}
\item ¿Cuál es la solución de la ecuación fraccionaria: \(\frac{x}{3} - 10 = 4\)?
\begin{enumerate}[label=\textbf{\Alph*)}]
\begin{multicols}{2}
  \item \(x = 18\)
  \item \(x = 30\)
  \item \(x = 42\)
  \item \(x = 48\)
\end{multicols}
\end{enumerate}
\vspace{0.15cm}
\item Determina el valor de la incógnita en: \(3x + 2x - 5 = 20\)
\begin{enumerate}[label=\textbf{\Alph*)}]
\begin{multicols}{2}
  \item \(x = 3\)
  \item \(x = 5\)
  \item \(x = 6\)
  \item \(x = 8\)
\end{multicols}
\end{enumerate}
\vspace{0.15cm}
\item Si se tiene la ecuación lineal \(3x = 24\), ¿cuánto vale \(x\)?
\begin{enumerate}[label=\textbf{\Alph*)}]
\begin{multicols}{2}
  \item \(x = 6\)
  \item \(x = 7\)
  \item \(x = 8\)
  \item \(x = 12\)
\end{multicols}
\end{enumerate}
\vspace{0.15cm}
\item ¿Cuál de las siguientes expresiones corresponde a la estructura de una función exponencial?
\begin{enumerate}[label=\textbf{\Alph*)}]
\begin{multicols}{2}
  \item \(f(x) = 10x + 5\)
  \item \(f(x) = 10^x\)
  \item \(f(x) = x^2\)
  \item \(f(x) = \sqrt{x}\)
\end{multicols}
\end{enumerate}
\vspace{0.15cm}
\item Dada la función exponencial \(f(x) = 5^x\), ¿cuál es el valor de \(f(4)\)?
\begin{enumerate}[label=\textbf{\Alph*)}]
\begin{multicols}{2}
  \item 20
  \item 125
  \item 625
  \item 3125
\end{multicols}
\end{enumerate}
\vspace{0.15cm}
\item En la función logarítmica: si \(\log_2(16) = x\), ¿cuál es el valor de \(x\)?
\begin{enumerate}[label=\textbf{\Alph*)}]
\begin{multicols}{2}
  \item 2
  \item 3
  \item 4
  \item 8
\end{multicols}
\end{enumerate}
\vspace{0.15cm}
\item ¿Qué valor numérico representa el logaritmo común en base 10 de 1000: \(\log_{10}(1000)\)?
\begin{enumerate}[label=\textbf{\Alph*)}]
\begin{multicols}{2}
  \item 2
  \item 3
  \item 10
  \item 100
\end{multicols}
\end{enumerate}
\vspace{0.15cm}
\item ¿Qué propiedad algebraica de los logaritmos corresponde a \(\log(\frac{a}{b})\)?
\begin{enumerate}[label=\textbf{\Alph*)}]
\begin{multicols}{2}
  \item \(\log(a) + \log(b)\)
  \item \(\log(a) - \log(b)\)
  \item \(\log(a) \times \log(b)\)
  \item \(\frac{\log(a)}{\log(b)}\)
\end{multicols}
\end{enumerate}
\vspace{0.15cm}
\item Dados los conjuntos \(A = \{1, 2\}\) y \(B = \{2, 3\}\), ¿es verdad que la unión \(A \cup B = \{1, 2, 3\}\)?
\begin{enumerate}[label=\textbf{\Alph*)}]
\begin{multicols}{2}
  \item Sí, es correcto
  \item No, debe ser \(\{2\}\)
  \item No, debe ser \(\{1, 2, 2, 3\}\)
  \item No, es el conjunto vacío
\end{multicols}
\end{enumerate}
\vspace{0.15cm}
\item Si \(A = \{1, 2, 3, 4, 5\}\) y \(B = \{5, 9, 13, 15, 19\}\), ¿cuál es la intersección \(A \cap B\)?
\begin{enumerate}[label=\textbf{\Alph*)}]
\begin{multicols}{2}
  \item \(\{1, 5, 19\}\)
  \item \(\{5\}\)
  \item \(\varnothing\) (conjunto vacío)
  \item \(\{1, 2, 3, 4, 5, 9, 13, 15, 19\}\)
\end{multicols}
\end{enumerate}
\vspace{0.15cm}
\item Sean los conjuntos \(A = \{a, b\}\) y \(B = \{b, c\}\). ¿Cuál es el resultado de la unión \(A \cup B\)?
\begin{enumerate}[label=\textbf{\Alph*)}]
\begin{multicols}{2}
  \item \(\{b\}\)
  \item \(\{a, b, c\}\)
  \item \(\{a, c\}\)
  \item \(\{a, b, b, c\}\)
\end{multicols}
\end{enumerate}
\vspace{0.15cm}
\item ¿Qué símbolo matemático universal se emplea para denotar la operación de Unión entre conjuntos?
\begin{enumerate}[label=\textbf{\Alph*)}]
\begin{multicols}{2}
  \item \(\cap\)
  \item \(\cup\)
  \item \(\subset\)
  \item \(\in\)
\end{multicols}
\end{enumerate}
\vspace{0.15cm}
\item Efectúa la suma de fracciones homogéneas: \(\frac{2}{26} + \frac{24}{26}\)
\begin{enumerate}[label=\textbf{\Alph*)}]
\begin{multicols}{2}
  \item \(\frac{26}{52}\)
  \item 1
  \item \(\frac{12}{13}\)
  \item 2
\end{multicols}
\end{enumerate}
\vspace{0.15cm}
\item Calcula el producto de las siguientes fracciones: \(\frac{5}{4} \times \frac{2}{5}\)
\begin{enumerate}[label=\textbf{\Alph*)}]
\begin{multicols}{2}
  \item \(\frac{7}{9}\)
  \item \(\frac{1}{2}\)
  \item \(\frac{10}{9}\)
  \item \(\frac{5}{8}\)
\end{multicols}
\end{enumerate}
\vspace{0.15cm}
\item Resuelve la sustracción de fracciones: \(\frac{5}{6} - \frac{1}{3}\)
\begin{enumerate}[label=\textbf{\Alph*)}]
\begin{multicols}{2}
  \item \(\frac{4}{3}\)
  \item \(\frac{1}{2}\)
  \item \(\frac{2}{3}\)
  \item \(\frac{1}{6}\)
\end{multicols}
\end{enumerate}
\vspace{0.15cm}
\item Efectúa la división de fracciones: \(\frac{3}{4} \div \frac{2}{3}\)
\begin{enumerate}[label=\textbf{\Alph*)}]
\begin{multicols}{2}
  \item \(\frac{6}{12}\)
  \item \(\frac{9}{8}\)
  \item \(\frac{8}{9}\)
  \item \(\frac{5}{7}\)
\end{multicols}
\end{enumerate}
\vspace{0.15cm}
\item Realiza la siguiente adición de números decimales: \(9.5 + 4.5 + 5.5 + 0.5\)
\begin{enumerate}[label=\textbf{\Alph*)}]
\begin{multicols}{2}
  \item 19.0
  \item 20.0
  \item 20.5
  \item 21.0
\end{multicols}
\end{enumerate}
\vspace{0.15cm}
\item ¿Cuál es el cociente de dividir \(5.25 \div 0.5\)?
\begin{enumerate}[label=\textbf{\Alph*)}]
\begin{multicols}{2}
  \item 2.625
  \item 10.5
  \item 10.05
  \item 25.0
\end{multicols}
\end{enumerate}
\vspace{0.15cm}
\item Calcula el resultado de la multiplicación: \(0.75 \times 8\)
\begin{enumerate}[label=\textbf{\Alph*)}]
\begin{multicols}{2}
  \item 5.5
  \item 6.0
  \item 6.5
  \item 7.0
\end{multicols}
\end{enumerate}
\vspace{0.15cm}
\item Resuelve la operación: \(5.2 \times 0.3\)
\begin{enumerate}[label=\textbf{\Alph*)}]
\begin{multicols}{2}
  \item 0.156
  \item 1.56
  \item 15.6
  \item 1.8
\end{multicols}
\end{enumerate}
\vspace{0.15cm}
\item Si un departamento salvadoreño tiene 348,880 habitantes de una población nacional de 6,359,223, ¿qué porcentaje representa aproximadamente?
\begin{enumerate}[label=\textbf{\Alph*)}]
\begin{multicols}{2}
  \item 3.2\%
  \item 5.5\%
  \item 8.0\%
  \item 12.5\%
\end{multicols}
\end{enumerate}
\vspace{0.15cm}
\item ¿Cuál es el área de una figura cuadrada cuyo lado mide exactamente 10 cm?
\begin{enumerate}[label=\textbf{\Alph*)}]
\begin{multicols}{2}
  \item 40 cm²
  \item 100 cm²
  \item 1000 cm²
  \item 20 cm²
\end{multicols}
\end{enumerate}
\vspace{0.15cm}
\item ¿Cuál es el perímetro total de un rectángulo de 10 cm de longitud de largo y 5 cm de ancho?
\begin{enumerate}[label=\textbf{\Alph*)}]
\begin{multicols}{2}
  \item 15 cm
  \item 30 cm
  \item 50 cm
  \item 60 cm
\end{multicols}
\end{enumerate}
\vspace{0.15cm}
\end{enumerate}\n
\vspace{0.4cm}
\begin{tcolorbox}[colback=gray!5!white,colframe=gray!60!black,title=\textbf{CLAVE DE RESPUESTAS -- MATEMÁTICAS}]
\begin{multicols}{5}
\begin{enumerate}[label=\textbf{\arabic*.}]
\setcounter{enumi}{205}
  \item \textbf{B}
  \item \textbf{B}
  \item \textbf{C}
  \item \textbf{B}
  \item \textbf{C}
  \item \textbf{B}
  \item \textbf{C}
  \item \textbf{C}
  \item \textbf{B}
  \item \textbf{B}
  \item \textbf{A}
  \item \textbf{B}
  \item \textbf{B}
  \item \textbf{B}
  \item \textbf{B}
  \item \textbf{B}
  \item \textbf{B}
  \item \textbf{B}
  \item \textbf{B}
  \item \textbf{B}
  \item \textbf{B}
  \item \textbf{B}
  \item \textbf{B}
  \item \textbf{B}
  \item \textbf{B}
\end{enumerate}
\end{multicols}
\end{tcolorbox}
\newpage\n
% ==============================================================================
% MATERIA 9: ESTADÍSTICA Y PROBABILIDAD
% ==============================================================================
\section{9. Estadística y Probabilidad (10 Preguntas)}

\begin{enumerate}[label=\textbf{\arabic*.}]
\setcounter{enumi}{227}
\item Si en una investigación estadística se encuesta a 10,000 personas de El Salvador para inferir hábitos nacionales, ese grupo seleccionado es:
\begin{enumerate}[label=\textbf{\Alph*)}]
\begin{multicols}{2}
  \item El censo universal
  \item Una muestra de la población
  \item Un parámetro absoluto
  \item Una constante matemática
\end{multicols}
\end{enumerate}
\vspace{0.15cm}
\item ¿Cómo se define técnicamente una "muestra representativa" en estadística?
\begin{enumerate}[label=\textbf{\Alph*)}]
\begin{multicols}{2}
  \item Toda la población completa
  \item Un subconjunto que reproduce fidedignamente las características esenciales de la población
  \item Un grupo de una sola familia
  \item Los datos erróneos de un estudio
\end{multicols}
\end{enumerate}
\vspace{0.15cm}
\item ¿A qué se refiere el término metodológico "sesgo de selección" en un muestreo?
\begin{enumerate}[label=\textbf{\Alph*)}]
\begin{multicols}{2}
  \item A un error de imprenta
  \item A una muestra viciada o no representativa que favorece indebidamente a un subgrupo
  \item Al cálculo de la media aritmética
  \item A la recolección exhaustiva de datos
\end{multicols}
\end{enumerate}
\vspace{0.15cm}
\item Al lanzar un dado común de seis caras balanceado, ¿cuál es la probabilidad teórica de obtener exactamente la cara 4?
\begin{enumerate}[label=\textbf{\Alph*)}]
\begin{multicols}{2}
  \item \(\frac{1}{4}\)
  \item \(\frac{1}{6}\)
  \item \(\frac{4}{6}\)
  \item \(\frac{1}{2}\)
\end{multicols}
\end{enumerate}
\vspace{0.15cm}
\item ¿Cuál es la probabilidad matemática de obtener "corona" (cara) al lanzar al aire una moneda justa?
\begin{enumerate}[label=\textbf{\Alph*)}]
\begin{multicols}{2}
  \item \(\frac{1}{4}\)
  \item \(\frac{1}{2}\) (50\%)
  \item \(\frac{1}{3}\)
  \item 1
\end{multicols}
\end{enumerate}
\vspace{0.15cm}
\item Si se arroja un dado común de 6 caras, ¿cuál es la probabilidad de que resulte un número impar?
\begin{enumerate}[label=\textbf{\Alph*)}]
\begin{multicols}{2}
  \item \(\frac{1}{6}\)
  \item \(\frac{1}{2}\) (3 de 6)
  \item \(\frac{2}{3}\)
  \item \(\frac{5}{6}\)
\end{multicols}
\end{enumerate}
\vspace{0.15cm}
\item ¿Es verdadero que un candado de seguridad con combinación de 3 dígitos (del 0 al 9 cada uno) admite 1,000 combinaciones posibles?
\begin{enumerate}[label=\textbf{\Alph*)}]
\begin{multicols}{2}
  \item Sí (de 000 a 999, es decir \(10 \times 10 \times 10 = 1000\))
  \item No, solo 30 combinaciones
  \item No, admite 300 combinaciones
  \item No, son infinitas
\end{multicols}
\end{enumerate}
\vspace{0.15cm}
\item Si una persona dispone de 3 camisas de diferente color y 2 pantalones distintos, ¿cuántos atuendos combinados puede vestir?
\begin{enumerate}[label=\textbf{\Alph*)}]
\begin{multicols}{2}
  \item 5 atuendos
  \item 6 atuendos
  \item 8 atuendos
  \item 9 atuendos
\end{multicols}
\end{enumerate}
\vspace{0.15cm}
\item ¿De cuántas formas distintas pueden ordenarse 3 libros diferentes en un estante lineal?
\begin{enumerate}[label=\textbf{\Alph*)}]
\begin{multicols}{2}
  \item 3 formas
  \item 6 formas
  \item 9 formas
  \item 12 formas
\end{multicols}
\end{enumerate}
\vspace{0.15cm}
\item ¿Qué postulado afirma que si un evento ocurre de \(m\) formas y otro sucesivo de \(n\) formas, ambos ocurren de \(m \times n\) formas?
\begin{enumerate}[label=\textbf{\Alph*)}]
\begin{multicols}{2}
  \item Principio aditivo
  \item Principio fundamental de multiplicación (conteo)
  \item Teorema de Pitágoras
  \item Ley de los grandes números
\end{multicols}
\end{enumerate}
\vspace{0.15cm}
\end{enumerate}\n
\vspace{0.4cm}
\begin{tcolorbox}[colback=gray!5!white,colframe=gray!60!black,title=\textbf{CLAVE DE RESPUESTAS -- ESTADÍSTICA Y PROBABILIDAD}]
\begin{multicols}{5}
\begin{enumerate}[label=\textbf{\arabic*.}]
\setcounter{enumi}{227}
  \item \textbf{B}
  \item \textbf{B}
  \item \textbf{B}
  \item \textbf{B}
  \item \textbf{B}
  \item \textbf{B}
  \item \textbf{A}
  \item \textbf{B}
  \item \textbf{B}
  \item \textbf{B}
\end{enumerate}
\end{multicols}
\end{tcolorbox}
\newpage\n
% ==============================================================================
% MATERIA 10: OFIMÁTICA E INTELIGENCIA ARTIFICIAL
% ==============================================================================
\section{10. Ofimática e Inteligencia Artificial (10 Preguntas)}

\begin{enumerate}[label=\textbf{\arabic*.}]
\setcounter{enumi}{240}
\item ¿Qué software ofimático líder permite crear, redactar, editar y dar formato profesional a documentos de texto?
\begin{enumerate}[label=\textbf{\Alph*)}]
\begin{multicols}{2}
  \item Microsoft Access
  \item Microsoft Word
  \item Microsoft Excel
  \item Microsoft Publisher
\end{multicols}
\end{enumerate}
\vspace{0.15cm}
\item ¿En qué pestaña de la cinta de opciones de Microsoft Word se localizan las herramientas de corrección ortográfica y gramatical?
\begin{enumerate}[label=\textbf{\Alph*)}]
\begin{multicols}{2}
  \item Pestaña 'Insertar'
  \item Pestaña 'Revisar'
  \item Pestaña 'Disposición'
  \item Pestaña 'Vista'
\end{multicols}
\end{enumerate}
\vspace{0.15cm}
\item ¿Cuál es el atajo de teclado para forzar un salto de página inmediato en Microsoft Word?
\begin{enumerate}[label=\textbf{\Alph*)}]
\begin{multicols}{2}
  \item Shift + Enter
  \item Ctrl + Enter
  \item Alt + Enter
  \item Ctrl + P
\end{multicols}
\end{enumerate}
\vspace{0.15cm}
\item ¿Qué tecla combinada con el ratón permite seleccionar fragmentos o bloques de texto no contiguos en un documento?
\begin{enumerate}[label=\textbf{\Alph*)}]
\begin{multicols}{2}
  \item Tab
  \item Ctrl
  \item Alt Gr
  \item Bloq Mayús
\end{multicols}
\end{enumerate}
\vspace{0.15cm}
\item ¿Cuál es el límite máximo predeterminado de diapositivas que puede contener una presentación de PowerPoint?
\begin{enumerate}[label=\textbf{\Alph*)}]
\begin{multicols}{2}
  \item Máximo 100 diapositivas
  \item No tiene un límite predeterminado (depende solo de la memoria del equipo)
  \item Exactamente 500
  \item Limitado a 50
\end{multicols}
\end{enumerate}
\vspace{0.15cm}
\item ¿En qué consiste la función "Vista del Moderador / Presentador" en Microsoft PowerPoint?
\begin{enumerate}[label=\textbf{\Alph*)}]
\begin{multicols}{2}
  \item En apagar la pantalla secundaria
  \item En un panel donde el expositor visualiza notas privadas, cronómetro y la siguiente diapositiva mientras el público ve solo la presentación
  \item En traducir el audio a texto automáticamente
  \item En bloquear el teclado
\end{multicols}
\end{enumerate}
\vspace{0.15cm}
\item ¿Qué herramienta permite configurar efectos de animación visual al pasar de una diapositiva a la siguiente en PowerPoint?
\begin{enumerate}[label=\textbf{\Alph*)}]
\begin{multicols}{2}
  \item Hipervínculos
  \item Transiciones
  \item Macros
  \item Formatos de celda
\end{multicols}
\end{enumerate}
\vspace{0.15cm}
\item En la versión en inglés de Microsoft Excel, ¿qué sintaxis suma correctamente los valores del rango A1 hasta A10?
\begin{enumerate}[label=\textbf{\Alph*)}]
\begin{multicols}{2}
  \item =ADD(A1:A10)
  \item =SUM(A1:A10)
  \item =TOTAL(A1:A10)
  \item =PLUS(A1:A10)
\end{multicols}
\end{enumerate}
\vspace{0.15cm}
\item En Microsoft Excel configurado en español, ¿qué fórmula suma los valores contenidos en las celdas desde A1 hasta A10?
\begin{enumerate}[label=\textbf{\Alph*)}]
\begin{multicols}{2}
  \item =ADICION(A1:A10)
  \item =SUMA(A1:A10)
  \item =SUMAR(A1..A10)
  \item =TOTAL(A1-A10)
\end{multicols}
\end{enumerate}
\vspace{0.15cm}
\item En Microsoft Excel en español, ¿qué función calcula la media aritmética de los valores del rango A1 a A5?
\begin{enumerate}[label=\textbf{\Alph*)}]
\begin{multicols}{2}
  \item =MEDIA(A1:A5)
  \item =PROMEDIO(A1:A5)
  \item =MEDIANA(A1:A5)
  \item =AVERAGE(A1:A5)
\end{multicols}
\end{enumerate}
\vspace{0.15cm}
\end{enumerate}\n
\vspace{0.4cm}
\begin{tcolorbox}[colback=gray!5!white,colframe=gray!60!black,title=\textbf{CLAVE DE RESPUESTAS -- OFIMÁTICA E INTELIGENCIA ARTIFICIAL}]
\begin{multicols}{5}
\begin{enumerate}[label=\textbf{\arabic*.}]
\setcounter{enumi}{240}
  \item \textbf{B}
  \item \textbf{B}
  \item \textbf{B}
  \item \textbf{B}
  \item \textbf{B}
  \item \textbf{B}
  \item \textbf{B}
  \item \textbf{B}
  \item \textbf{B}
  \item \textbf{B}
\end{enumerate}
\end{multicols}
\end{tcolorbox}
\newpage\n
\end{document}